\documentclass[11pt]{article}
\usepackage[T1]{fontenc}
\usepackage{lmodern}
\usepackage{amsmath,amssymb,amsthm,mathtools}
\usepackage[margin=1in]{geometry}
\usepackage{microtype,booktabs,longtable}
\usepackage{xurl}
\usepackage[hidelinks]{hyperref}
\newtheorem{theorem}{Theorem}[section]
\newtheorem{lemma}[theorem]{Lemma}
\newtheorem{proposition}[theorem]{Proposition}
\newtheorem{corollary}[theorem]{Corollary}
\theoremstyle{definition}
\newtheorem{definition}[theorem]{Definition}
\theoremstyle{remark}
\newtheorem{remark}[theorem]{Remark}
\newcommand{\Z}{\mathbb Z}
\newcommand{\R}{\mathbb R}
\newcommand{\Q}{\mathbb Q}
\newcommand{\T}{\mathbb R/\mathbb Z}
\newcommand{\F}{\mathbb F}
\newcommand{\E}{\mathbb E}
\newcommand{\Prob}{\mathbb P}
\newcommand{\ceil}[1]{\left\lceil #1\right\rceil}
\newcommand{\floor}[1]{\left\lfloor #1\right\rfloor}
\newcommand{\set}[1]{\left\{#1\right\}}
\DeclareMathOperator{\Tr}{Tr}
\DeclareMathOperator{\rank}{rank}
\DeclareMathOperator{\sat}{sat}
\DeclareMathOperator{\dist}{dist}
\allowdisplaybreaks
\title{Affine compression and largest Sidon subsets\\\large Research report and proof audit}
\author{Research session with Codex}
\date{3 October 2026 --- final session report}
\begin{document}
\maketitle

\begin{abstract}
For every $n$-element set of integers, we prove the existence of a Sidon
subset of size
\[
\left(\frac{2}{3\sqrt3}+o(1)\right)\sqrt n
=\bigl(0.3849001794\ldots+o(1)\bigr)\sqrt n.
\]
The proof strengthens the compression estimate supplied in the task by
adding a common random affine shift and using exact two-point averaging.
The resulting injective Freiman map has a domain of size at least
$n/2-n(n-1)/(8m)$ in a cyclic group of arbitrary order $m$.
We also prove a weighted version with uniform point-inclusion probabilities.
This does not resolve Erd\H{o}s's conjectural constant $1$.
The argument has independent internal audits; publication novelty has not
been certified. In particular, affine randomization already appears in
Ruzsa's 1995 work, whose consequences require correcting the supplied
claim about the previous best benchmark.
\end{abstract}
\tableofcontents

\section{Result and status}

Throughout, a Sidon set means a set for which
\[
b_1+b_2=b_3+b_4
\quad\Longrightarrow\quad
\{b_1,b_2\}=\{b_3,b_4\}
\]
as unordered multisets. Repetitions are allowed; in particular a
nonconstant three-term arithmetic progression is forbidden. Write $s(A)$
for the largest cardinality of a Sidon subset of $A$, and put
\[
H(n)=\min_{A\subset\Z,\ |A|=n}s(A).
\]

\begin{theorem}[Universal lower bound]\label{thm:main}
For every prime power $q$, put $m=q^2+q+1$. Every $n$-element set
$A\subset\Z$ contains a Sidon subset of cardinality at least
\begin{equation}\label{eq:finite-bound}
\ceil{\frac{q+1}{m}
\ceil{\frac n2-\frac{n(n-1)}{8m}}}.
\end{equation}
A negative right side can of course be replaced by zero, or by one if
$n\ge1$. Consequently
\begin{equation}\label{eq:main}
H(n)\ge\left(\frac{2}{3\sqrt3}+o(1)\right)\sqrt n.
\end{equation}
The error term depends only on $n$, not on the diameter or additive
structure of $A$. The same lower bound holds for finite subsets of
$\R^N$, uniformly in $N$.
\end{theorem}

Theorem~\ref{thm:main} is a proved result of this session. It improves the
specific Bailleul--Riblet bound supplied in the task by a factor of two.
It is not a proof or a disproof of the conjecture $H(n)\sim\sqrt n$.
All statements described below as obstructions concern specified methods;
none supplies an asymptotic counterexample to that conjecture.

The finite arithmetic checks accompanying this report corroborate the
constructions and conventions. No universal or asymptotic assertion in
the main proof depends on a numerical experiment.

\section{Literature audit and quantitative comparison}

Bailleul--Riblet~\cite{BR}, Lemma~2.3, applies to every finite
$A\subset\Z$ and every positive integer $m$, with no diameter or density
restriction. It gives an injective Freiman $2$-homomorphism on a subset of
size at least $n/2-n^2/(2m)$. Their proof first obtains the slightly
stronger subtraction $\binom n2/m$. Its Freiman implication preserves
existing equalities; it is not asserted to reflect all equalities. Taking
a ceiling costs nothing. Theorem~2.2 optimizes at $m\sim3n$.
Their Lemma~2.5 uses Singer's covering; translating one Singer set gives
the slightly stronger exact intersection average $(q+1)|C|/m$ used here.

The starting paper's theorem is verified, but its description as the best
available universal constant cannot be adopted without qualification.
Ruzsa~\cite{RuzsaII}, Theorem~4.4, gives $r'(N)\ge4r(N)/(S^2+4)$ for
translation-invariant integral systems with maximum row $\ell^1$ norm
$S$, allowing specified coincidence patterns. With the row
$(1,1,-1,-1)$ and all nontrivial Sidon coincidence patterns, $S=4$,
$r'(N)=H(N)$, and $r(N)$ is the interval Sidon extremum. Thus that
theorem implies the stronger prior consequence
$H(N)\ge(1/5+o(1))\sqrt N$. A direct check is given below.
Ruzsa already uses a common affine shift and occupied intervals; the
randomization mechanism is not claimed as new here.

\begin{center}
\begin{tabular}{@{}lll@{}}
\toprule
Source or deduction & Constant multiplying $\sqrt n$ & Status\\
\midrule
Bailleul--Riblet, Theorem 2.2 & $1/(3\sqrt3)=0.192450\ldots$ & Verified source statement\\
Ruzsa, Theorem 4.4, Sidon specialization & $1/5=0.2$ & Verified older consequence\\
Optimized interval-template argument & $\sqrt2/(3\sqrt3)=0.272166\ldots$ & Derived in this audit\\
Theorem~\ref{thm:main} & $2/(3\sqrt3)=0.384900\ldots$ & Proved here\\
Erd\H{o}s's conjecture & $1$ & Unresolved here\\
\bottomrule
\end{tabular}
\end{center}

A bounded search through 3 October 2026 found no subsequent primary source
stating the universal $0.384900\ldots$ constant. Huber's August 2026
paper~\cite{Huber} still quotes the Bailleul--Riblet benchmark for arbitrary
integer sets; its other constants concern a particular multinomial
family. Pach--Zakharov~\cite{PZ}, Proposition~2, is a close predecessor for
random floor compression. Results for intervals, weak Sidon ambient
sets, squares, or other restricted families are not substitutes for a
universal statement about $H(n)$. This bounded search is not a novelty
certificate. The complete text of Ruzsa's 2006 paper on additive and
multiplicative Sidon sets~\cite{Ruzsa2006} was not obtained, which is
an additional limit of this priority check.

One unresolved bibliographical point is recorded rather than silently
used: Green's exercise sheet~\cite{Green}, Question~7, states a large
subset admitting an injective Freiman $s$-homomorphism into $[n]$, but
provides no proof. Its literal $n/(2s)$ bound was not independently
reconstructed in this session. No theorem in this report relies on it.

\subsection{Direct verification of the older transfer and its optimization}

Let $D\subseteq\{1,\ldots,M\}$ be nonempty and Sidon, with $d=|D|$, and put
\[
J_y=\left(\frac{y-1/4}{2M},\frac{y+1/4}{2M}\right)
\quad(y\in D),\qquad \ell=|J_y|=\frac1{4M}.
\]
Project an arbitrary $n$-element integer set by $a\mapsto\{a\theta+t\}$,
and keep one source point for each occupied $J_y$. This always gives a
Sidon set. Indeed, in a source equality write each circle image as
$(y_i+e_i)/(2M)$ with $|e_i|<1/4$. The integer
$y_1+y_2-y_3-y_4$ differs from a multiple of $2M$ by less than one,
so it equals that multiple. Its absolute value is at most $2(M-1)$,
so it is zero. The Sidon property of $D$ and the chosen distinct
representatives then trivialize the original equality.

Let $X$ be total hits, $Z$ the sum of squared occupancies and $Y$ the
number of occupied arcs. The two-point averaging proved in
Lemma~\ref{lem:pair} gives
\[
\E X=nd\ell,\qquad \E Z=nd\ell+n(n-1)d\ell^2.
\]
Pointwise $X\le\sqrt{YZ}$, and Cauchy--Schwarz on the parameter space
therefore implies
\[
\E Y\ge\frac{(\E X)^2}{\E Z}
=\frac{nd\ell}{1+(n-1)\ell}.
\]
For $M=n$ this is at least $d/5$, independently verifying the Sidon
specialization of Ruzsa's theorem, including repeated summands.

The interval construction $d\ge(1+o(1))\sqrt M$ used here follows
already from Lemma~\ref{lem:singer}: choose the largest prime
$q\le\sqrt M-1$, so $q=(1+o(1))\sqrt M$ and $q^2+q+1\le M$,
and take integer representatives of its Singer set. The prime number
theorem implies that every fixed relative interval below a sufficiently
large real number contains a prime, which proves the asserted ratio.
An integer equality
implies an equality modulo the Singer modulus, so the representatives
are Sidon as integers.

The elementary occupied-cell inequality used in
Theorem~\ref{thm:compression} gives the better estimate
\[
\E Y\ge d\left(\frac n{4M}-\frac{n(n-1)}{32M^2}\right).
\]
For $M\sim\beta n$ its coefficient of $\sqrt n$ is
$1/(4\sqrt\beta)-1/(32\beta^{3/2})$.
Its derivative is $(3-8\beta)/(64\beta^{5/2})$, giving the maximum
$\sqrt2/(3\sqrt3)$ at $\beta=3/8$.
This optimized consequence is a calculation in this audit, not an
explicit constant attributed to Ruzsa's paper. The cyclic target below
removes the factor lost in fitting the interval template without wraparound.

\section{Exact affine compression}

A Freiman $2$-homomorphism $\phi:C\to G$ satisfies
$a+b=c+d\Rightarrow\phi(a)+\phi(b)=\phi(c)+\phi(d)$ for all entries
in $C$. Injectivity is a separate condition.

\begin{lemma}[Two-point averaging]\label{lem:pair}
Let $a\ne b$ be integers, and let $\theta,t$ be independent uniform
points of $[0,1)$. The pair
\[
\bigl(\{a\theta+t\},\{b\theta+t\}\bigr)
\]
is uniform on $[0,1)^2$.
\end{lemma}
\begin{proof}
For each fixed $\theta$, $X=\{a\theta+t\}$ is uniform, so $X$ is
independent of $\theta$. The variable $W=\{(b-a)\theta\}$ is uniform:
multiplication by the nonzero integer $b-a$ preserves Lebesgue measure
on the circle, as can be checked by splitting $[0,1)$ into $|b-a|$
equal intervals. Hence $X,W$ are independent uniform circle points.
The second displayed coordinate is $X+W$ modulo one, which, conditional
on $X$, is uniform. This proves the assertion.
\end{proof}

\begin{theorem}[Affine compression replacement]\label{thm:compression}
Let $A\subset\Z$ have cardinality $n$, let $m\ge1$ be an integer, and
let $0<\delta\le1/2$. There are $C\subseteq A$ and an injective Freiman
$2$-homomorphism $C\to\Z/m\Z$ with
\begin{equation}\label{eq:compression}
|C|\ge\ceil{\delta n-\frac{\delta^2 n(n-1)}{2m}}.
\end{equation}
In particular, $|C|\ge\ceil{n/2-n(n-1)/(8m)}$ is always possible.
\end{theorem}
\begin{proof}
Choose $\theta,t$ as in Lemma~\ref{lem:pair}. For every $a\in A$ write
\[
m(a\theta+t)=L_a+e_a,\qquad L_a\in\Z,\quad0\le e_a<1,
\]
and define
\[
B=\{a\in A:0\le e_a<\delta\},\qquad
\phi(a)=L_a\pmod m\quad(a\in B).
\]
Suppose $a+b=c+d$ in $B$, allowing repeated entries. The common affine
term appears twice on each side, and therefore cancels. Consequently
\[
L_a+L_b-L_c-L_d=-(e_a+e_b-e_c-e_d).
\]
The left side is an integer. Both residual sums lie in $[0,2\delta)$,
so their difference lies strictly between $-1$ and $1$, including at
$\delta=1/2$. The integer is zero. Thus $\phi$ preserves the required
equalities, first over the integers and hence modulo $m$.

Every point is retained with probability $\delta$, including $a=0$.
For $j=0,\ldots,m-1$ put
\[
I_j=\left[\frac j m,\frac{j+\delta}{m}\right).
\]
A pair of distinct points $a,b$ is retained in the same $\phi$-fiber
exactly when their circle images lie in the same $I_j$.
Lemma~\ref{lem:pair} therefore gives the exact probability
\begin{equation}\label{eq:pair-collision}
\Prob(a,b\in B,\ \phi(a)=\phi(b))
=\sum_{j=0}^{m-1}|I_j|^2=\frac{\delta^2}{m}.
\end{equation}

Let $r_j$ be the number of retained points in fiber $j$. Keep one point
from each nonempty fiber to form $C$. This makes the restriction
injective, and does not destroy preservation of equalities. For every
integer $r\ge0$,
\[
\mathbf1_{r>0}\ge r-\binom r2:
\]
for $r\ge1$ the difference is $(r-1)(r-2)/2\ge0$, and $r=0$ is
immediate. Hence
\[
|C|\ge |B|-\sum_j\binom{r_j}{2}.
\]
Taking expectations and using~\eqref{eq:pair-collision} gives
\[
\E|C|\ge\delta n-\binom n2\frac{\delta^2}{m}.
\]
The integer-valued variable $|C|$ takes finitely many values, so some
choice is at least the ceiling of the right side. This proves the theorem.
No estimate involves the diameter of $A$.
\end{proof}

\begin{remark}[The shift must have the correct scale]
Here $t$ is uniform on $[0,1)$ in $\floor{m(a\theta+t)}$.
Equivalently, the additive shift in $\floor{am\theta+\tau}$ is uniform
on $[0,m)$. Taking $\tau$ uniform only on $[0,1)$ does not justify the
two-point Haar calculation as stated. Also, after the affine shift one
must use the balanced four-term identity above, rather than asserting
additivity of $a\mapsto\phi(a)$ on arbitrary sums.
\end{remark}

\section{Singer extraction and the optimized constant}

\begin{lemma}[Pullback through an injective forward map]\label{lem:pullback}
If $\phi:C\to\Z/m\Z$ is injective and is a Freiman $2$-homomorphism,
and $D\subseteq\Z/m\Z$ is Sidon, then $C$ contains a Sidon subset of
size at least $\ceil{|C||D|/m}$.
\end{lemma}
\begin{proof}
Average $|\phi(C)\cap(D+u)|$ over all $u\in\Z/m\Z$. Each point of
$\phi(C)$ is counted $|D|$ times, so some intersection has size at least
$|C||D|/m$. Pull it back to $S\subset C$. Any original equality in $S$
maps, in the forward direction, to an equality in $D+u$. Its two image
multisets agree by the Sidon property. Injectivity then makes the
original multisets agree. This includes equal summands and uses no
reflection of additive equalities.
\end{proof}

For completeness, the finite-geometric input is proved here in its
usual trace form, rather than relying on a covering assertion.

\begin{lemma}[Singer]\label{lem:singer}
For every prime power $q$, the cyclic group of order $m=q^2+q+1$
contains a Sidon set of size $q+1$.
\end{lemma}
\begin{proof}
Let $K=\F_{q^3}$ and $F=\F_q$. The quotient $K^*/F^*$ is cyclic of
order $m$. Define $\Tr(x)=x+x^q+x^{q^2}$. This is a nonzero
$F$-linear map $K\to F$: the displayed nonzero polynomial has degree
$q^2<q^3$, so cannot vanish on all of $K$. Its kernel has $q^2$ elements.
The nonzero elements of this kernel are a union of multiplicative
$F^*$-classes. Thus
\[
D=\{xF^*:x\ne0,\ \Tr(x)=0\}
\]
has cardinality $(q^2-1)/(q-1)=q+1$.

For $\gamma F^*\ne F^*$, the forms $x\mapsto\Tr(x)$ and
$x\mapsto\Tr(\gamma^{-1}x)$ are independent. Indeed, if the second
were $c$ times the first, then
$\Tr((\gamma^{-1}-c)x)=0$ for every $x$. If $\gamma^{-1}-c\ne0$,
multiplication by this element is a bijection of $K$, contradicting
nonvanishing of the trace. Thus $\gamma^{-1}=c\in F$, a contradiction.
The two kernels consequently intersect in a one-dimensional $F$-space.
Their nonzero intersection gives exactly one multiplicative $F^*$-class.
It follows that every nonidentity element of $K^*/F^*$ occurs exactly
once as an ordered difference of two distinct members of $D$, in
additive notation for the cyclic quotient.

Finally, if $x+y=z+w$ with $x,y,z,w\in D$, then either $x=z$, giving
$y=w$, or $x-z=w-y$ is a nonzero difference. Its uniqueness gives
$(x,z)=(w,y)$. This is exactly the Sidon assertion, including the
diagonal cases.
\end{proof}

\begin{proof}[Proof of Theorem~\ref{thm:main} over $\Z$]
Apply Theorem~\ref{thm:compression} with $\delta=1/2$, followed by
Lemmas~\ref{lem:singer} and~\ref{lem:pullback}. This proves the exact
finite bound~\eqref{eq:finite-bound}.

For a fixed $\beta>1/4$, choose primes
$q=(1+o(1))\sqrt{\beta n}$. Such choices exist by the prime number
theorem: for every fixed $\eta>0$, the interval $[x,(1+\eta)x]$
contains a prime for all sufficiently large $x$; consequently the least
prime above $x$ has ratio tending to one. Thus
\[
m=q^2+q+1=(\beta+o(1))n,\qquad
\frac{q+1}{m}=\frac{1+o(1)}{\sqrt{\beta n}}.
\]
The resulting coefficient of $\sqrt n$ is
\[
f(\beta)=\frac{1/2-1/(8\beta)}{\sqrt\beta}
=\frac{4\beta-1}{8\beta^{3/2}}.
\]
Its derivative is $(3-4\beta)/(16\beta^{5/2})$, so its maximum occurs
at $\beta=3/4$ and equals $2/(3\sqrt3)$.
At this choice the retained fraction tends to $1/3$.
Ceilings affect only bounded additive terms before extraction. All
choices depend on $n$ alone, giving the claimed uniformity.
\end{proof}

\begin{lemma}[Transfer from real configurations]\label{lem:real}
Every finite subset of $\R^N$ admits an injective, additive-relation
preserving map into $\Z$.
\end{lemma}
\begin{proof}
Take a basis over $\Q$ of the finite-dimensional rational span of the
given points. Express all points in this basis and clear a common
denominator; this gives integer vectors $v_1,\ldots,v_n\in\Z^d$,
with all original additive equalities preserved and distinct points
remaining distinct. For each nonzero difference $v_i-v_j$, the polynomial
$\sum_{r=1}^d(v_{i,r}-v_{j,r})X^{r-1}$ is nonzero and has finitely many
integer roots. Choose an integer $M$ outside the union of these finite
root sets. The linear map $v\mapsto\sum_r v_rM^{r-1}$ is then injective
on the points and preserves every additive equality. Pulling back a
Sidon subset of its image proves the remaining assertion of
Theorem~\ref{thm:main}. There is no dimension-dependent loss.
\end{proof}

\section{Weighted and balanced-equation forms}

\begin{theorem}[Uniform marginal guarantee]\label{thm:weighted}
Under the hypotheses of Theorem~\ref{thm:compression}, let
$D\subseteq\Z/m\Z$ be Sidon. There is a probability distribution on
Sidon subsets $S\subseteq A$ for which every $a\in A$ satisfies
\begin{equation}\label{eq:marginal}
\Prob(a\in S)\ge\frac{|D|}{m}
\left(\delta-\frac{(n-1)\delta^2}{2m}\right).
\end{equation}
In particular, for arbitrary nonnegative weights $w_a$, some Sidon
subset has at least this fraction of the total weight. With the
optimized Singer parameters, this fraction is
$(2/(3\sqrt3)+o(1))/\sqrt n$, uniformly over all weight assignments.
\end{theorem}
\begin{proof}
Use the affine parameters from the compression proof and assign the
points independent continuous random priorities. Keep the lowest-priority
point in each retained fiber. If a retained point $a$ is discarded,
some $b\ne a$ occupies its retained fiber and precedes it. For each
competitor this has probability $\delta^2/(2m)$, by
\eqref{eq:pair-collision} and independence of priorities. A union bound
therefore gives
\[
\Prob(a\in C)\ge\delta-\frac{(n-1)\delta^2}{2m}.
\]
Choose independently a uniform translate $D+u$ and take its pullback
inside $C$. Conditional on $C,\phi$, a fixed surviving point belongs
to the pullback with probability $|D|/m$. The pullback is always Sidon
by Lemma~\ref{lem:pullback}, proving~\eqref{eq:marginal}. Multiply by
$w_a$, sum, and take expectations. The distribution itself is independent
of the weights.
\end{proof}

Dividing the finitely many output-subset probabilities by a positive
lower bound in~\eqref{eq:marginal} gives a fractional cover of $A$ by
Sidon subsets: the weights of sets containing any point sum to at least
one. Its total weight is at most
$(3\sqrt3/2+o(1))\sqrt n$.

\begin{proposition}[Balanced equations]\label{prop:balanced}
Let $h\ge2$ be an integer. The same construction gives an injective
Freiman $h$-homomorphism from a subset of every $n$-element integer set
into $\Z/m\Z$, of cardinality at least
\[
\ceil{\frac n h-\frac{n(n-1)}{2h^2m}}.
\]
More generally it preserves every balanced integral equation whose
positive coefficient sum is at most $h$, with its variables in the
retained subset.
\end{proposition}
\begin{proof}
Set $\delta=1/h$. In a balanced equation the affine shifts cancel.
The positive and negative weighted sums of fractional remainders each
belong to $[0,1)$, since the positive and negative coefficient sums
are equal and at most $h$. Their difference is an integer strictly
between $-1$ and $1$. The collision calculation is unchanged and gives
the stated cardinality. Repeated variables simply contribute their
integer multiplicities to this calculation.
\end{proof}

The application justifying this replacement is
Theorem~\ref{thm:main}, with all its losses already accounted for.
The broader balanced-equation wording is not being presented as a
separate improvement without an extremal target construction.

\subsection{A quantitative surplus for progression-rich sets}

\begin{theorem}[Progression-sensitive compression]\label{thm:ap-surplus}
For a finite integer set $A$, let $T(A)$ count its nontrivial three-term
arithmetic progressions, each counted once as $a<b<c$ with $a+c=2b$.
For $n=|A|$, $m\ge1$, and $0<\delta\le1/2$, there is a subset
$C\subset A$ and an injective Freiman $2$-homomorphism
$C\to\Z/m\Z$ with
\[
|C|\ge\left\lceil\delta n-
\frac{\delta^2}{2m}\bigl(n(n-1)-T(A)\bigr)\right\rceil.
\]
Consequently, for any fixed $0\le\tau\le1/4$, families satisfying
$T(A)\ge(\tau+o(1))n^2$ have Sidon subsets of size at least
\[
\left(\frac{2}{3\sqrt{3(1-\tau)}}+o(1)\right)\sqrt n.
\]
\end{theorem}
\begin{proof}
Use the affine phases and retained cells of
Theorem~\ref{thm:compression}. Let $k_j$ count points in retained cell
$j$, and let $O$ count occupied cells. The exact identity
\[
O=\sum_j k_j-\sum_j\binom{k_j}{2}+\sum_j z(k_j),
\qquad z(0)=0,\quad z(k)=\frac{(k-1)(k-2)}2\ (k\ge1),
\]
exhibits the slack in pair deletion. A set of $k$ distinct integers
has at most
\[
\sum_{i=0}^{k-1}\min(i,k-1-i)
=\left\lfloor\frac{(k-1)^2}{4}\right\rfloor
\]
three-term progressions: each endpoint on one side of a specified
midpoint determines the endpoint on its other side. For $k\ge3$ this
bound is at most $z(k)$; for $k\le2$ both relevant quantities vanish.
Thus $\sum_jz(k_j)$ is at least the number of source progressions
contained in individual retained cells.

For a fixed progression $a,a+d,a+2d$ with $d>0$, put
$X=\{a\theta+t\}$ and $V=\{d\theta\}$. These are independent
uniform circle points, by the same argument as Lemma~\ref{lem:pair}.
The endpoint phases are $X,Y=X+2V$. Conditional on $X,Y$, the variable
$V$ is uniform among its two possible roots; hence the middle phase
$X+V$ is uniform among the two roots of $2z=X+Y$.
This argument works for every positive integer $d$.
For any arc of length $\epsilon\le1/2$, when both endpoints lie in
it their lifted midpoint is in the arc and its half-turn translate
is outside, apart from null boundaries. Therefore the three-hit
probability for this arc is exactly $\epsilon^2/2$.
Taking $\epsilon=\delta/m$ and summing the $m$ retained cells gives
probability $\delta^2/(2m)$ for each progression. Pairwise independence
now yields
\[
\E O\ge\delta n-\binom n2\frac{\delta^2}{m}
+T(A)\frac{\delta^2}{2m}.
\]
Choose parameters attaining at least the ceiling of this average and
keep one point per occupied cell. The Freiman implication, including
diagonals, follows from the same residual-integer calculation as in
Theorem~\ref{thm:compression}.

For the Sidon assertion take $\delta=1/2$ and Singer moduli
$m\sim3(1-\tau)n/4$, available by prime asymptotics. Translating and
averaging the Singer set gives the coefficient
\[
\max_{\beta>0}
\frac{1/2-(1-\tau)/(8\beta)}{\sqrt\beta}
=\frac{2}{3\sqrt{3(1-\tau)}}.
\]
All estimates depend on $n,T(A),m$ alone, independently of diameter.
\end{proof}

This is a strict improvement over the universal bound when $\tau>0$,
but does not settle the progression-sparse branch. For intervals,
$T(A)=\lfloor(n-1)^2/4\rfloor$, and the coefficient furnished by this
argument is $4/9$; stronger interval results were already known.
Boundedly many separated long interval clusters also have a positive
progression density. The claim is a uniform bound for the stated
class, not a new universal coefficient or a claimed novelty result.

\begin{proposition}[Local progression-sensitive marginals]\label{prop:weighted-ap}
Let $A\subset\Z$ have $n$ elements, let $m\ge1$ be an integer,
and let $T_a(A)$ count the nontrivial three-term progressions in $A$
containing $a$. For $0<\delta\le1/2$, there is a weight-independent
distribution on subsets $C$ admitting injective Freiman
$2$-homomorphisms $C\to\Z/m\Z$, such that every $a\in A$ satisfies
\[
\Prob(a\in C)\ge\delta-
\frac{\delta^2}{2m}\left(n-1-\frac13T_a(A)\right).
\]
If $m=q^2+q+1$, a subsequent independent uniform Singer translate
gives a Sidon subset $S$ with each of these lower bounds multiplied
by $(q+1)/m$. Thus for nonnegative weights $w_a$, writing
$W=\sum_a w_a$, some such $S$ satisfies
\[
\sum_{a\in S}w_a\ge\frac{q+1}{m}
\left[\delta W-\frac{\delta^2}{2m}
\left((n-1)W-\frac13\sum_a w_aT_a(A)\right)\right].
\]
The distribution of $S$ does not depend on the weights.
\end{proposition}
\begin{proof}
We construct a representative distribution for each retained cell $F$
of size $r$. Let $d_a$ count its progressions through $a$. Define
\[
L_a=\max\left\{0,\frac{3-r}{2}+\frac{d_a}{3}\right\}.
\]
We claim $\sum_aL_a\le1$. For $r=1,2$ the demands are respectively
$1$ and $1/2$. For $r\ge3$, translate $a$ to zero. Every progression
through it is a doubling edge $x,2x$ or an opposite pair $x,-x$ among
the other differences. Doubling edges form a forest within each sign;
opposite pairs form a matching. If $a$ has increasing rank $i$, starting
at zero, this gives
\[
d_a\le r-3+\min(i,r-1-i)\quad(0<i<r-1),\qquad
d_a\le r-2\quad(i=0,r-1).
\]
For $r=3$ every demand is at most $1/3$; for $r=4$ every demand is at
most $1/6$. For $r\ge5$ the endpoint demands vanish, and interior
demands satisfy
\[
L_a\le\max\left\{0,\frac{3-r+2\min(i,r-1-i)}6\right\}.
\]
For odd $r$, only the median can have positive demand, at most $1/3$.
For even $r$, only the two central ranks can, each at most $1/6$.
The claim follows. Choose one representative with probabilities
$p_a=L_a+(1-\sum_{b\in F}L_b)/r$. This depends on $F$, not on the
weights, and gives $p_a\ge1-(r-1)/2+d_a/3$.

Multiply the conditional inequality by the indicator that $a$ is
retained and average. Each other point shares its retained cell with
probability $\delta^2/m$. Each progression through $a$ lies in one
retained cell with probability $\delta^2/(2m)$, as proved in
Theorem~\ref{thm:ap-surplus}. This proves the marginal estimate.
An independent Singer translate multiplies all marginals by
$(q+1)/m$. Multiplication by $w_a$ and summation prove the weighted
claim. All source equalities and diagonals are preserved by the same
compression map already constructed. The coefficient $1/3$ is optimal
for this pointwise comparison: in a cell consisting of a three-term
progression, $r=3$ and all three degrees are one, so a coefficient $c$
would require three probabilities at least $c$, forcing $c\le1/3$.
\end{proof}

For example, if a fixed $0\le\eta\le3/4$ satisfies
$T_a(A)\ge(\eta+o(1))n$ uniformly in $a$, optimizing
at $\delta=1/2$ and $m\sim3(1-\eta/3)n/4$ gives a distribution of
Sidon subsets with every marginal at least
\[
\left(\frac{2}{3\sqrt{3(1-\eta/3)}}+o(1)\right)n^{-1/2}.
\]
This requires the stated local progression hypothesis. Summing the
marginals counts each progression three times and recovers the full
$T(A)$ correction of Theorem~\ref{thm:ap-surplus}.

\section{Relation to the no-unique-sum manuscript}

The local manuscript's Proposition \texttt{prop:compression} assumes an
inclusion-minimal no-unique-sum set in $\Z/p\Z$ of minimum size $n(p)$.
It produces a primitive sublattice of the collision lattice, with free
quotient rank $O(\sqrt{n(p)}+n(p)/\log p)$ and polynomially bounded
coordinates in one quotient basis. Its proof obtains departing
relations from small block rectangles, contracts forests using unit
pivots, and retains residual relations. This report checked those
hypotheses and the mechanism; it does not use its quantitative conclusion
as an input or certify the full source manuscript theorem.

The original hypothesis cannot hold for a nonempty finite integer set:
$2\min A$ has a unique unordered representation. A low-dimensional free
partial quotient, by itself, also does not establish injectivity into
a small finite group.

Here is the precise character-space interpretation of the replacement.
Let $L\le\Z^n$ be generated by all actual collision rows
$e_i+e_j-e_k-e_l$, including diagonal pairs, and let
$Q=\Z^n/\sat(L)$. This quotient is free abelian. Both augmentation
$e_i\mapsto1$ and evaluation $e_i\mapsto a_i$ annihilate $\sat(L)$,
because their target $\Z$ is torsion-free. Thus the point classes
$v_i$ are distinct, and every pair $v_i,v_j$, $i\ne j$, is rationally
independent: augmentation gives $r+s=0$ in any relation
$rv_i+sv_j=0$, and evaluation gives $r(a_i-a_j)=0$.

A uniform character $Q\to\T$ consequently has independent uniform
values on every two distinct point classes. To check this explicitly,
choose a basis of $Q$. The two corresponding rows of integers have rank
two, so the real linear map to $\R^2$ is onto; its torus map is onto and
pushes Haar measure to Haar measure. All original additive relations
remain satisfied modulo one. If their chosen real lifts satisfy
$x_a+x_b-x_c-x_d=k\in\Z$, multiplying by $m$ shows that the
rounded-label difference equals $mk$ minus the residual-sum difference.
The latter is an integer in $(-1,1)$, hence zero. Thus half-window
rounding still preserves the equality modulo $m$, and pair deletion
gives Theorem~\ref{thm:compression}. In fact, the rank-two
family $a_i\theta+t$ already provides everything needed for that theorem.

All residual equations must be imposed. A primitive partial quotient
$\Z^n/P$ with $P\subset L$ allows characters that violate the remaining
rows, so cannot replace $Q$ without further work. Saturation is harmless
for this integer application because augmentation and separating integer
evaluation descend through it. It would generally destroy the modular
rank defect sought in the original no-unique-sum problem, so this is an
application-specific replacement, not a torsion-preserving upgrade of
the original method.

% Fragment for sidon_research_report.tex. No independent preamble.
% Uses the main report's theorem environments and \Z,\F,\sat,\ceil macros.
\section{Exact homomorphisms: a useful criterion and a sharp structural warning}
\label{sec:content-hashing}

This section gives a conditional route to the larger coefficient
$2\sqrt2/(3\sqrt3)=0.544331\ldots$, followed by a proved obstruction to
making its hypotheses universal through the full relation group alone.
The criterion is an application-specific replacement for an elimination
lemma, not a claim that the supplied no-unique-sum manuscript establishes
it. No novelty claim is made for the restricted-class bound.

\subsection{Contents in the complete Freiman quotient}

Let $A=\{a_1,\ldots,a_n\}\subset\Z$, and let $L_A\le\Z^n$ be generated
by every actual relation row
\[
e_i+e_j-e_k-e_l\qquad(a_i+a_j=a_k+a_l),
\]
including rows with repeated indices. Set
\[
Q_A=\Z^n/\sat(L_A),\qquad
\sat(L_A)=\{x\in\Z^n:rx\in L_A\text{ for some }r\ge1\},
\]
and write $v_i$ for the class of $e_i$. This is a free abelian group.
The maps $e_i\mapsto1$ and $e_i\mapsto a_i$ annihilate the saturation:
they annihilate $L_A$ and take values in the torsion-free group $\Z$.
In particular the $v_i$ are distinct.
For nonzero $w\in Q_A$, define its \emph{content} $c(w)$ to be the
largest positive integer $c$ with $w\in cQ_A$. In any integral basis it
is the gcd of the coordinates, so it is independent of the basis.

\begin{lemma}[Exact content hashing]\label{lem:content-hash}
For every integer $m\ge1$, put $c_{ij}=c(v_i-v_j)$. There exist
$C\subseteq A$ and an injective Freiman $2$-homomorphism
$C\to\Z/m\Z$ such that
\begin{equation}\label{eq:content-bound}
|C|\ge\ceil{n-\frac1m\sum_{i<j}\gcd(c_{ij},m)}.
\end{equation}
\end{lemma}
\begin{proof}
Send the members of a basis of $Q_A$ independently and uniformly into
$\Z/m\Z$, obtaining a random homomorphism $h$. For an integer coordinate
vector $(w_1,\ldots,w_r)$, the image $h(w)$ is uniform on the subgroup
it generates modulo $m$. That subgroup consists of the multiples of
$g=\gcd(w_1,\ldots,w_r,m)=\gcd(c(w),m)$ and has size $m/g$.
Thus
\[
\Prob(h(v_i)=h(v_j))=\frac{\gcd(c_{ij},m)}m.
\]
The expected number of colliding unordered pairs is the sum in
\eqref{eq:content-bound} divided by $m$. Keeping one member of each
nonempty fiber discards $k-1\le\binom k2$ points from a fiber of size
$k$. Some $h$ therefore admits the required injective restriction.
Every source equality is preserved because $h$ annihilates $L_A$.
The ceiling follows from integrality.
\end{proof}

If every $c_{ij}$ is coprime to $m$, this gives
$|C|\ge\ceil{n-n(n-1)/(2m)}$. For Singer moduli $m\sim\beta n$,
translation averaging then gives coefficient
\[
\frac{1-1/(2\beta)}{\sqrt\beta},
\]
whose maximum, at $\beta=3/2$, is $2\sqrt2/(3\sqrt3)$.
More generally, for fixed $\epsilon\ge0$, if suitable Singer moduli satisfy
\[
\sum_{i<j}(\gcd(c_{ij},m)-1)\le(\epsilon+o(1))n^2,
\qquad m\sim\frac{3(1+2\epsilon)}2n,
\]
the same calculation gives coefficient
\[
\frac{2\sqrt2}{3\sqrt3\sqrt{1+2\epsilon}}.
\]
The existence of these moduli with the indicated content control is a
hypothesis, not a conclusion. Merely bounding the \emph{number} of
nonprimitive pairs by $\epsilon n^2$ is insufficient: an individual
bad pair can collide with probability one.

\subsection{A finite-field formulation retaining residual arithmetic}

For an odd prime $p$, use instead the full modular relation space
\[
V_p=\F_p^n/\operatorname{span}_{\F_p}
\{e_i+e_j-e_k-e_l:a_i+a_j=a_k+a_l\},
\]
with point classes $u_i$. Here the integral relation lattice is
\emph{not} saturated before reduction. Thus possible $p$-torsion and
modular rank drops are retained.

\begin{proposition}[Modular separation criterion]\label{prop:modular-separation}
If the classes $u_1,\ldots,u_n$ are distinct, then
\begin{equation}\label{eq:field-bound}
s(A)\ge\ceil{\frac np-\frac{n(n-1)}{2p^3}}.
\end{equation}
Consequently, if this hypothesis holds along primes
$p\sim\sqrt{3n/2}$, then
\[
s(A)\ge\left(\frac{2\sqrt2}{3\sqrt3}+o(1)\right)\sqrt n.
\]
\end{proposition}
\begin{proof}
A uniformly random linear map $h:V_p\to\F_p^2$ sends each nonzero
$u_i-u_j$ uniformly into $\F_p^2$. The pair collision probability is
therefore $p^{-2}$. Keeping one point per image, as in
Lemma~\ref{lem:content-hash}, gives an injectively mapped subset $C$
with $|C|\ge n-n(n-1)/(2p^2)$, preserving every source equality.

The parabola
\[
S=\{(x,x^2):x\in\F_p\}\subset\F_p^2
\]
is a strict Sidon set of size $p$. Equality of two pair sums gives
$x+y=z+w$ and $x^2+y^2=z^2+w^2$. Since $2$ is invertible, it also gives
$xy=zw$. The two unordered pairs, with multiplicity, are the roots of
the same monic quadratic and hence agree. This includes repeated
summands. Averaging all translates of $S$ gives a translate meeting
$h(C)$ in at least $|C|/p$ points. Its preimage is Sidon by the forward
Freiman implication and injectivity. This proves
\eqref{eq:field-bound}. Writing $p\sim t\sqrt n$, the coefficient is
$t^{-1}-(2t^3)^{-1}$, maximized at $t=\sqrt{3/2}$.
\end{proof}

\begin{corollary}[Box models]\label{cor:box-model}
Suppose $A$ has an injective Freiman $2$-homomorphism into
$\{0,\ldots,K\}^d$, with arbitrary $d$.
For each odd prime $p>K$, the bound~\eqref{eq:field-bound} holds.
For each fixed $\eta>0$,
the coefficient $2\sqrt2/(3\sqrt3)$ holds uniformly in $d$ whenever
\[
K\le\bigl(\sqrt{3/2}-\eta\bigr)\sqrt n.
\]
\end{corollary}
\begin{proof}
Reducing the coordinate map modulo $p$ preserves all relations and
remains injective: a nonzero coordinate difference has absolute value
at most $K<p$. This map factors through $V_p$, so its point classes
are distinct. Apply Proposition~\ref{prop:modular-separation}.
The prime number theorem supplies primes $p\sim\sqrt{3n/2}$, which
exceed $K$ for all sufficiently large $n$ under the displayed condition.
\end{proof}

For example, any subset of
$\{\sum_jd_jB^j:0\le d_j\le K\}$, with $B>2K$, has the required digit
map: two-term digit sums have no carries. The dimension and the integer
diameter are unrestricted. This is a concrete class to which the
criterion applies, not a literature-optimality claim for that class.
If integral elimination instead gives point coordinates in
$[-H,H]^d$, translating coordinates makes $K=2H$ available. All
residual equations must still be satisfied; a bounded coordinate list
in a partial quotient is not sufficient.

If the criterion holds only on a subset of size
$k=(\alpha+o(1))n$, at primes $p\sim\sqrt{3k/2}$, its coefficient for
$\sqrt n$ is
$\frac{2\sqrt2}{3\sqrt3}\sqrt\alpha$. It improves the universal
$2/(3\sqrt3)$ coefficient precisely when $\alpha>1/2$.
No decomposition guaranteeing this retention, or treating every
complementary case, is proved here.

\subsection{Small sets of anchors can force simultaneous collapse}

The next result explains why the complete relation group cannot simply
be computed once and then projected to various small target groups.
It is a method obstruction, not a counterexample to the Sidon conjecture.

\begin{theorem}[Simultaneous rigidifier]\label{thm:rigidifier}
For each $n$, let $\mathcal M_n$ be a finite nonempty family of positive
integer moduli, and suppose
\[
L_n=\operatorname{lcm}(\mathcal M_n),\qquad \log L_n=o(n).
\]
For all sufficiently large $n$ there is an integer set $A_n$ of
\emph{exactly} $n$ points with the following properties.
Its complete relation lattice is saturated, and its full Freiman group
is freely generated by two elements $z,w$, with
\[
v_a=z+aw\qquad(a\in A_n).
\]
The set contains a progression
\[
P_n=\{L_n,2L_n,\ldots,N_nL_n\},\qquad N_n=n-o(n).
\]
For every $m\in\mathcal M_n$ and every homomorphism
$h:Q_{A_n}\to\Z/m\Z$, all points of $P_n$ have the same image.
Every injective restriction of such an $h$ therefore has size
$O(1+\log L_n)=o(n)$.
\end{theorem}
\begin{proof}
Write $L=L_n$. Begin with $0,1$ and read the binary expansion of $L$
after its leading digit. From the current integer $x$, include $2x$;
if the next digit is one, also include $2x+1$, which becomes the new
current integer. Otherwise $2x$ becomes current. Let $T$ be all points
so produced, including $0,1$. Then $L\in T$, every other point lies
in $[0,L)$, and
\[
t:=|T|\le2\floor{\log_2 L}+2.
\]
For $L=1$ simply use $T=\{0,1\}$. Prescribe
$N=n-t+1>0$ and set
\[
A=T\cup\{L,2L,\ldots,NL\}.
\]
The intersection is exactly $\{L\}$, so $|A|=n$ and
$N=n-o(n)$.

Put $z=e_0$ and $w=e_1-e_0$. The actual chain relations are
\[
e_x+e_x=e_{2x}+e_0,
\qquad e_{2x}+e_1=e_{2x+1}+e_0
\]
when the indicated steps occur. They successively force
$e_t=z+tw$ for all $t\in T$. For $k=2,\ldots,N$, the actual relation
\[
e_{(k-1)L}+e_L=e_{kL}+e_0
\]
then forces $e_{kL}=z+kLw$.

There is exactly one displayed relation for each generator other than
$e_0,e_1$. In its construction order each has coefficient $-1$ on its
new generator and involves only earlier generators otherwise.
Eliminating these unit pivots leaves the free group on $z,w$, with no
residual relation. Equivalently, the displayed rows generate the
kernel of the surjection
\[
\Z^A\longrightarrow\Z^2,\qquad e_a\longmapsto(1,a).
\]
Every actual additive relation lies in this kernel. Conversely, every
displayed row is an actual relation. Thus their subgroup is the
\emph{entire} relation lattice, already saturated, and the claimed
full quotient follows.

For $m\mid L$, every homomorphism into $\Z/m\Z$ satisfies
$h(v_{kL})=h(z)+kLh(w)=h(z)$. An injective restriction keeps at most
one progression point and at most $t-1$ other points. Its size is at
most $t=O(1+\log L)=o(n)$, as required.
\end{proof}

The same proof applies to any abelian target whose exponent divides
$L_n$, including $\F_p^r$ for every $r$ whenever $p\mid L_n$.
Nevertheless, deleting only $t-1=o(n)$ anchors leaves $P_n$.
Its \emph{rebuilt} Freiman model is $kL_n\mapsto k$; it embeds into
$\Z/m\Z$ for $m\ge N_n$ and contains a Sidon subset of size
$(1+o(1))\sqrt{N_n}$ by the interval construction.
There is therefore no Sidon obstruction. The failure is specific to
restricting homomorphisms of the old full quotient after subset selection.

\begin{corollary}[All linear-sized Singer targets]\label{cor:singer-rigidifier}
Fix $C>0$. There are $n$-element integer sets on which every exact
homomorphism from the full Freiman group into every Singer cyclic group
of order at most $Cn$ has injective restrictions of size $o(n)$.
\end{corollary}
\begin{proof}
Use the family $q^2+q+1$ for all integers
$2\le q\le\floor{\sqrt{Cn}}$, which contains every desired prime-power
Singer modulus. Its lcm divides the product of these
$O(\sqrt n)$ integers, each $O(n)$, so its logarithm is
$O(\sqrt n\log n)=o(n)$. Apply Theorem~\ref{thm:rigidifier}.
No assertion about prime gaps is needed.
\end{proof}

\subsection{Every modulus in a prescribed short window}

\begin{lemma}[An interval lcm bound]\label{lem:interval-lcm}
For integers $M\ge k\ge1$,
\[
\operatorname{lcm}(M-k+1,\ldots,M)
\ \mid\ \binom Mk\operatorname{lcm}(1,\ldots,k).
\]
In particular, if $M=O(n)$ and $k=o(n)$, the logarithm of the left
side is $o(n)$.
\end{lemma}
\begin{proof}
Fix a prime $p$. Let $a$ be the largest $p$-adic valuation among the
$k$ integers in the interval, and put $b=\floor{\log_p k}$.
For each $j$, that interval contains at least $\floor{k/p^j}$ multiples
of $p^j$. For $b<j\le a$ it contains at least one, while
$\floor{k/p^j}=0$. The factorial formula therefore gives
$v_p\binom Mk\ge\max(a-b,0)$. Adding the valuation $b$ of
$\operatorname{lcm}(1,\ldots,k)$ proves the divisibility.

We supply the needed elementary estimate
$\log\operatorname{lcm}(1,\ldots,k)=O(k)$.
Let $\psi(x)=\sum_{p^a\le x}\log p$. There is at most one power of
any fixed prime in $(r,2r]$. Its corresponding term in the factorial
valuation of $\binom{2r}r$ is one, and all other terms are nonnegative.
Consequently
$\psi(2r)-\psi(r)\le\log\binom{2r}r\le2r\log2$.
Summing over powers of two gives $\psi(k)\le4k\log2$.
Finally,
\[
\log\operatorname{lcm}(M-k+1,\ldots,M)
\le k\log(eM/k)+O(k)=o(n),
\]
since $(k/n)\log(n/k)\to0$.
\end{proof}

\begin{corollary}[Prescribed short windows]\label{cor:window-rigidifier}
Fix $\beta>0$ and a prescribed sequence $w_n\ge1$ with $w_n=o(n)$.
There exist sets $A_n\subset\Z$ of exactly $n$ points such that every
exact homomorphism from their full Freiman group into $\Z/m\Z$, for
\emph{every} positive integer $m$ with
$|m-\beta n|\le w_n$, has injective restrictions of size $o(n)$.
The same holds for a fixed finite union of such prescribed windows.
\end{corollary}
\begin{proof}
The moduli form an interval of $k\le2w_n+1=o(n)$ consecutive integers
with largest element $M=O(n)$. Its lcm has logarithm $o(n)$ by
Lemma~\ref{lem:interval-lcm}. For a fixed finite union, the lcm divides
the product of the separate lcms and has the same little-$o$ property.
Apply Theorem~\ref{thm:rigidifier}, which prescribes final size $n$
\emph{before} choosing the progression length. This avoids any mismatch
between the center of the modulus window and the constructed set size.
\end{proof}

The window in this corollary is prescribed. It does not assert that a
single sequence defeats every possible unspecified rate in $m\sim\beta n$.
By contrast, Corollary~\ref{cor:singer-rigidifier} covers every Singer
modulus up to $Cn$, without a rate qualification.
Neither statement blocks maps from a quotient rebuilt on a subset, or
the affine rounding construction proved earlier.

The accompanying exact-integer scripts
\path{research/sidon/scripts/rigidifier_certificate.py} and
\path{research/sidon/scripts/short_window_certificate.py}
check finite instances. The first uses $N=64$, moduli $7,13,31$,
$L=2821$, and $n=80$, and integrally reduces $2041$ generating
source relations against the unit-pivot basis. The second prescribes
$n=256$ and every modulus $188\le m\le196$; it has $180$ progression
points and verifies $16438$ relation generators, as well as all
$20100$ interval-lcm cases $1\le k\le M\le200$.
Their JSON outputs are in \texttt{research/sidon/results/}.
The proofs above, including all asymptotic claims, are independent of
these computations.

% Fragment for sidon_research_report.tex. No independent preamble.
\section{A robust obstruction even after rebuilding the quotient}
\label{sec:robust-cluster-obstruction}

The preceding obstruction concerns homomorphisms inherited from the full
set. The following construction is stronger in a different direction:
it allows arbitrary new Freiman maps after deletion, and arbitrary finite
abelian target groups. Its elementary proof is recorded here without a
claim of publication novelty.

\begin{theorem}[Robust finite-model obstruction]
\label{thm:robust-cluster}
Fix integers $r\ge2$ and $N\ge6$, and put
\[
B_r=\{0,\ldots,r-1\}\cup\{r,2r,\ldots,r^2\},\qquad
A_{r,N}=\{3Nb+x:b\in B_r,\ 0\le x<N\}.
\]
Thus $|A_{r,N}|=2rN$. Suppose $C\subseteq A_{r,N}$,
$|A_{r,N}\setminus C|=K<N$.
Every Freiman $2$-homomorphism $f:C\to G$ into an abelian group has
the form
\[
f(3Nb+x)=\gamma+b\alpha+x\beta\qquad(3Nb+x\in C)
\]
for some $\gamma,\alpha,\beta\in G$. If $f$ is injective and $G$ is
finite, then
\begin{equation}\label{eq:robust-cluster-order}
|G|\ge(r^2+1)(N-K).
\end{equation}
For $K=0$, equality is attained by a cyclic target group.
The condition $K<N$ for the affine-rigidity conclusion is sharp:
at $K=N$ that conclusion can fail, even for $G=\Z$.
\end{theorem}

\begin{lemma}[Affine rigidity of a dense interval]
\label{lem:dense-interval-affine}
Let $S\subseteq\{0,\ldots,N-1\}$ omit $k$ points, where $N\ge3k+3$.
Every Freiman $2$-homomorphism $f:S\to G$ into any abelian group is
of the form $f(x)=t+x\beta$ on $S$. The threshold is sharp for each
$k\ge1$: at $N=3k+2$ this conclusion can fail.
\end{lemma}
\begin{proof}
If $d\ge0$ occurs as a difference of two points of $S$, the increment
\[
g(d)=f(x+d)-f(x)
\]
is independent of the choice of $x$: the equality
$(x+d)+y=x+(y+d)$ equates any two such increments.
When $k=0$, consecutive increments immediately prove the assertion.
Otherwise let $M=\{0,\ldots,N-1\}\setminus S$, and let $g_0\ge1$
be the length of a longest consecutive run in $M$. For
$0\le d\le g_0$, a starting integer $x$ fails one of the conditions
$x,x+d,x+d+1\in S$ only if it belongs to
\[
M\cup(M-d)\cup(M-d-1).
\]
The last two sets overlap in at least $g_0-1$ points, using two
unit translates of a longest missing run. Their union with $M$
therefore has cardinality at most $3k-g_0+1$.
There are $N-d-1\ge N-g_0-1\ge3k-g_0+2$ candidate starts with
$0\le x$ and $x+d+1<N$, so a valid triple exists.
It gives $g(d+1)=g(d)+g(1)$. Induction yields
$g(d)=d\beta$ for all $0\le d\le g_0+1$, where $\beta=g(1)$.
Every consecutive gap in $S$ is at most $g_0+1$.
Telescoping along these gaps gives
$f(y)-f(x)=(y-x)\beta$ for arbitrary $x,y\in S$.
Choosing a base point gives the asserted affine formula.
No division in $G$ was used.

For sharpness, let $k\ge1$, $N=3k+2$, and
$S=\{0,\ldots,k\}\cup\{2k+1,\ldots,3k+1\}$.
Its pair sums with zero, one, or two summands in the second block
lie respectively in the disjoint intervals
$[0,2k]$, $[2k+1,4k+1]$, and $[4k+2,6k+2]$.
The map equal to zero on the first block and one on the second
is therefore a Freiman $2$-homomorphism into $\Z$.
It is not affine: its values at $0$ and $1$ force both affine
parameters to vanish, contradicting its value at $2k+1$.
\end{proof}

\begin{proof}[Proof of Theorem~\ref{thm:robust-cluster}]
The map $(b,x)\mapsto3Nb+x$ is a Freiman $2$-isomorphism from
$B_r\times\{0,\ldots,N-1\}$ onto $A_{r,N}$:
in any two-sum equality the vertical discrepancy has absolute value
at most $2N-2<3N$, so both coordinates must agree separately.
We therefore use the product coordinates.

Write $S_b=\{x:(b,x)\in C\}$ and $k_b=N-|S_b|$, so
$\sum_b k_b=K$. Choose a least-deleted fiber $S_c$, and put
$h=k_c=\min_b k_b$. Its map is affine,
$f(c,x)=t_c+x\beta$. Indeed, for $h=0$ this follows from consecutive
increments on the full interval. For $h=1$ the dense-interval lemma
applies because $N\ge6$. For $h\ge2$, the inequalities
\[
N\ge K+1\ge2rh+1\ge4h+1\ge3h+3
\]
again permit that lemma.

Consider consecutive elements $x<y$ of any other fiber $S_b$, and
put $d=y-x$. The $d-1$ omitted positions between them and the omissions
from all other fibers give
\[
K\ge d-1+(2r-1)h.
\]
If $h>0$, it follows that $N-d\ge(2r-1)h>2h$. Among the $N-d$
pairs at difference $d$ in the full interval, at most $2h$ are spoiled
in $S_c$, so some $z,z+d\in S_c$ survive. If $h=0$, the full base
fiber supplies such a pair for every $d<N$ directly. The rectangle
relation
\[
(b,y)+(c,z)=(b,x)+(c,z+d)
\]
therefore gives $f(b,y)-f(b,x)=d\beta$. Telescoping along consecutive
elements shows $f(b,x)=t_b+x\beta$ throughout every fiber.
Single-point fibers satisfy this formula as well, and no fiber is empty
because $K<N$.

For every relation $b_1+b_2=b_3+b_4$ in $B_r$, the intersection of
the corresponding distinct fibers omits at most $K<N$ positions.
Evaluating at a common $x$ shows
$t_{b_1}+t_{b_2}=t_{b_3}+t_{b_4}$. The interval
$\{0,\ldots,r\}\subset B_r$ and the relations
$(j-1)+1=j+0$ for $2\le j\le r$ give
$t_j=t_0+j(t_1-t_0)$. The relations
$((j-1)r)+r=jr+0$ for $2\le j\le r$ give the same formula for
all remaining points of $B_r$. Set $\gamma=t_0$ and
$\alpha=t_1-t_0$. This proves the asserted affine formula.

Every integer $h$ with $1\le h\le r^2$ is $jr-i$ for some
$1\le j\le r$ and $0\le i\le r-1$. Hence
$B_r-B_r=\{-r^2,\ldots,r^2\}$.
For $0<|h|\le r^2$, choose $b,c\in B_r$ with $b-c=h$.
For every integer $t$ with $|t|<N-K$, there are $N-|t|>K$
pairs $x,y\in\{0,\ldots,N-1\}$ satisfying $x-y=t$.
At most $k_b+k_c\le K$ pairs have $x\notin S_b$ or $y\notin S_c$.
Thus some pair survives. Injectivity of $f$ implies
\[
h\alpha+t\beta\ne0.
\]
For $h=0$ and $0<|t|<N-K$, choose a fiber with
$k_b\le K/(2r)$. At most $2k_b\le K$ pairs in that fiber at
difference $t$ are spoiled, so injectivity similarly gives
$t\beta\ne0$.
It follows that
\[
\{0,\ldots,r^2\}\times\{0,\ldots,N-K-1\}\longrightarrow G,
\qquad (h,t)\longmapsto h\alpha+t\beta
\]
is injective: every nonzero difference of two inputs is one of the
forbidden relations just established. This proves
\eqref{eq:robust-cluster-order}.

When $K=0$, the map
$(b,x)\mapsto b+(r^2+1)x$ modulo $(r^2+1)N$ preserves source
equalities and is injective, since $0\le b\le r^2$ and the indicated
integer representatives lie in $[0,(r^2+1)N-1]$.

For sharpness of rigidity, delete the $N$ points in row $b=0$.
The map $g$ in Proposition~\ref{prop:cluster-unfolding} preserves all
horizontal two-sum relations, so $f(b,x)=g(b)$ is a Freiman map into
$\Z$. It is not affine in $(b,x)$: the vertical slope would be zero,
and $g(1)=0$, $g(2)=1$ would force $g(b)=b-1$, whereas
$g(2r)=r\ne2r-1$. This concerns rigidity, not a claim that the
target-order lower bound must be sharp at every deletion size.
\end{proof}

\begin{corollary}\label{cor:no-near-lossless-model}
For every fixed $M>0$, there is a sequence of integer sets $A$ with
$|A|\to\infty$ such that no subset of size $|A|-o(|A|)$ admits an
injective Freiman $2$-homomorphism into any finite abelian group of
order at most $M|A|$.
\end{corollary}
\begin{proof}
Choose a fixed $r\ge2$ with $(r^2+1)/(2r)>M$ and let $N\to\infty$.
For $n=2rN$ and $K=o(n)=o(N)$, the hypotheses $N\ge6$ and $K<N$ hold
eventually, and \eqref{eq:robust-cluster-order} gives
$|G|\ge((r^2+1)/(2r)-o(1))n>Mn$.
\end{proof}

This obstruction has only three affine parameters and integral
unit-coefficient recurrences. Even so, rebuilding the quotient after
deleting a negligible fraction of the points does not produce a finite
injective model of uniformly bounded linear size.
It narrows the secondary compression objective substantially.
It does \emph{not} contradict Erd\H{o}s's conjecture: no asymptotic
Sidon upper bound below $\sqrt{|A_{r,N}|}$ is proved here.

\begin{proposition}[Deleting the anchor row unfolds the product]
\label{prop:cluster-unfolding}
For the family in Theorem~\ref{thm:robust-cluster},
\[
s(A_{r,N})\ge(1-o(1))\sqrt{(2r-1)N}\qquad(N\to\infty,\ r\text{ fixed}).
\]
Its lower coefficient relative to $\sqrt{|A_{r,N}|}$ is therefore
$\sqrt{1-1/(2r)}$. More precisely, deleting the row with horizontal
coordinate zero leaves a set admitting a bijective Freiman
$2$-homomorphism onto an interval of length $(2r-1)N$.
\end{proposition}
\begin{proof}
On $B_r\setminus\{0\}$ define
\[
g(b)=b-1\quad(1\le b\le r),\qquad
g(jr)=r+j-2\quad(2\le j\le r).
\]
This bijects the horizontal set onto $\{0,\ldots,2r-2\}$ and preserves
all two-sum equalities. To verify the latter assertion, call
$\{1,\ldots,r\}$ the small part and $\{2r,\ldots,r^2\}$ the large part.
A small--small sum is at most $2r$, whereas any sum involving a large
point is greater than $2r$; within the small part $g$ is affine.
A mixed sum whose small point is less than $r$ is not divisible by $r$
and uniquely determines its two summands, by reduction modulo $r$.
Every remaining sum involves only multiples of $r$; on the entire
progression $r,2r,\ldots,r^2$, the formula is $g(jr)=r+j-2$,
so these equalities are preserved as well. This includes repeated
summands and exhausts all possible equalities.

It follows that $(b,x)\mapsto Ng(b)+x$ is a bijective Freiman
$2$-homomorphism from
$(B_r\setminus\{0\})\times\{0,\ldots,N-1\}$ onto
$\{0,\ldots,(2r-1)N-1\}$. Pull back an interval Sidon set of size
$(1-o(1))\sqrt{(2r-1)N}$, available from Singer's construction and
prime asymptotics. The forward implication and injectivity suffice.
\end{proof}

The robust obstruction therefore distinguishes deletion of $o(N)$ points
from deletion of the $N$ points in one anchor row. Its increasingly large
finite-model obstruction as $r$ grows comes with actual Sidon subsets
whose lower coefficient tends to $1$, rather than with evidence for a
counterexample to the conjecture.

The following upper bound is already a special case of
Riblet's Theorem~3.1(ii)~\cite{RibletIntervals}; we include a
self-contained proof for these products.
\begin{proposition}[Known upper bound for the obstructing family]
\label{prop:cluster-sidon-range}
For each fixed $r\ge2$, as $N\to\infty$,
\[
s(A_{r,N})\le(1+o(1))\sqrt{2rN}.
\]
\end{proposition}
\begin{proof}
More generally let $B$ be any fixed set of
$b$ integers and let $S\subset B\times\{0,\ldots,N-1\}$ be Sidon.
Write $S_c=\{x:(c,x)\in S\}$, $k_c=|S_c|$, and $s=|S|$.
Every positive vertical difference occurs at most once among all
the fibers. In fact two occurrences give an equality of two source
sums, and the Sidon condition makes the ordered difference pairs
identical. Consequently
$\sum_c\binom{k_c}{2}\le N-1$, so $s=O_b(\sqrt N)$ by
$s^2/b\le\sum_c k_c^2\le2N-2+s$.

Put $u=\floor{N^{3/4}}$ and let $R_c(y)$ count representations
$y=x+t$ with $x\in S_c$ and $1\le t\le u$. Cauchy--Schwarz,
applied within each fiber, gives
\[
\frac{u^2}{N+u}\sum_c k_c^2
\le\sum_{c,y}R_c(y)^2
\le su+2\sum_{d=1}^{u-1}(u-d)
=su+u(u-1).
\]
The second inequality uses the uniqueness, across all fibers,
of each positive difference $d$. Since $s/u=o(1)$ and $u=o(N)$,
it follows that $\sum_c k_c^2\le(1+o(1))N$ and hence
$s^2\le b\sum_c k_c^2\le(1+o(1))bN$.
Take $B=B_r$ and $b=2r$.
\end{proof}

\section{An exact gain from doubling suitable fibers}
This finite strengthening is proved, but no uniform asymptotic gain from it
has been established.  Let $m=q^2+q+1$, and let $D\subset\mathbb Z/m\mathbb Z$
be a Singer difference set.  Suppose $\phi:A_0\to\mathbb Z/m\mathbb Z$ is a
Freiman $2$-homomorphism, where $A_0$ is a finite set of integers; injectivity
is not assumed.  Write $O$ for its number of occupied fibers.  In each fiber
of size at least two, choose one pair and assign that fiber the pair's
positive difference.  Let $N_d$ count fibers assigned difference $d$.

\begin{proposition}\label{prop:fiber-matching}
Under these hypotheses, $A_0$ has a Sidon subset of size at least
\[
 \left\lceil\frac{(q+1)O}{m}
 +\sum_{d:N_d>0}\frac{(q+1)^2N_d}{m(q+N_d)}\right\rceil.
\]
\end{proposition}
\begin{proof}
For a translate $D+u$, retain one point from each occupied selected fiber.
For each assigned difference occurring in at least one selected fiber,
choose one such fiber and replace its retained point by the assigned pair.
Each fiber has only one assignment, so these choices never double a fiber
twice.  The resulting set is Sidon.  Indeed, applying $\phi$ to a putative
sum equality makes the unordered multisets of fibers equal.  If the equality
lies in one fiber, a set of at most two integers is Sidon.  Otherwise it
involves two different fibers, and a nontrivial equality forces their
positive internal differences to agree, contrary to the construction.
This argument includes repeated summands.

Take $u$ uniformly in $\mathbb Z/m\mathbb Z$.  A specified fiber is selected
with probability $(q+1)/m$.  Two distinct specified fibers are both selected
with probability $1/m$: their nonzero difference has exactly one ordered
representation in the Singer difference set.  Let $Y_d$ be the number of
selected fibers assigned $d$.  Thus
\[
 \mathbb EY_d=\frac{(q+1)N_d}{m},\qquad
 \mathbb EY_d^2=\frac{N_d(q+N_d)}m.
\]
For $N_d>0$, Cauchy--Schwarz gives
\[
 \mathbb P(Y_d>0)\ge
 \frac{(\mathbb EY_d)^2}{\mathbb EY_d^2}
 =\frac{(q+1)^2N_d}{m(q+N_d)}.
\]
The constructed Sidon set has one point per occupied selected fiber and
one additional point for every $d$ with $Y_d>0$.  Averaging its integer size
proves the result, including the ceiling.
\end{proof}

The additional summands have order $\min\{N_d/\sqrt m,1\}$.  One may average
the proposition over the random affine maps in the main theorem, with any
deterministic rule for assigning pairs.  What remains unproved is a useful
uniform lower bound for this gain.  In particular, rare large fibers prevent
pair counts alone from guaranteeing many multiply occupied fibers, and the following bound limits the average gain to $O(1)$ when there
are only $O(n)$ distinct positive differences. For a half-window affine
map, two points in one retained fiber with positive difference $d$
satisfy $\operatorname{dist}(d\theta,\Z)<1/(2m)$. For each nonzero
integer $d$ this event has probability $1/m$. The number of additional
points in the construction is at most the number of such resonant
differences. Thus its expected gain is at most $|\Delta^+(A)|/m$,
where $\Delta^+(A)=\{b-a:a,b\in A,\ a<b\}$. In particular this is
$O(1)$ for intervals at $m\sim\beta n$.

\section{Obstructions to immediate further improvements}

\subsection{Half-window rounding and target spacing}

For every $\delta>1/2$ the same floor map can fail even on a three-term
progression. Choose $0<\varepsilon<\min(2\delta-1,1)$, $m\ge5$, and
\[
A=\{0,1,2\},\qquad mt=\frac12-\frac\varepsilon2,
\qquad m\theta=\frac32+\frac\varepsilon2.
\]
All three fractional remainders are below $\delta$, but the distinct
floor labels are $0,2,3$. They violate the image of $0+2=1+1$ modulo
$m$. Injectivity does not fix this error.

More generally, if a measurable residual set $R\subset[0,1)$ is
universally carry-free in the sense that reduction modulo one is
injective on $R+R$, then $|R|\le1/2$. For a compact $K\subset R$,
the sets $\min(K)+K$ and $\max(K)+K$ overlap in at most one point,
so $|K+K|\ge2|K|$. Injectivity modulo one gives $|K+K|\le1$.
Approximating $R$ from inside by compact subsets proves the assertion.
For any such window of measure $\delta$, pair collisions still have
probability $\delta^2/m$; changing its shape does not change these
two moments.

An attempted repair with a full window requires a stronger target.
If a Sidon set $D\subset\Z/m\Z$ also has no adjacent elements in
$D+D$, then
\[
|D|(|D|-1)+1=|D-D|\le\floor{m/2}.
\]
Indeed, adjacent differences $a-b,c-d$ would give adjacent sums
$a+d,c+b$, and an adjacency-free subset of a cycle has at most
$\floor{m/2}$ vertices. Thus the stronger target has asymptotic size
at most $\sqrt{m/2}$. With a full rounding window, optimizing the same
pair-deletion estimate and this target density again gives at most
$2/(3\sqrt3)$. This only limits that specified combination of methods.

\subsection{A sharp limit for arbitrary interval templates}

The preceding limitation extends to unequal windows around arbitrary
circle points. This concerns the specified pair-deletion guarantee;
it does not upper-bound the actual occupied-cell count.

\begin{proposition}[Interval-template limit]\label{prop:template-limit}
Let $J_1,\ldots,J_s$ be disjoint positive-length arcs in $\T$, of
lengths $\ell_1,\ldots,\ell_s$. Suppose they have the following
universal label property: if $x_r\in J_{i_r}$ for $1\le r\le4$ and
$x_1+x_2=x_3+x_4$ in $\T$, then the unordered multisets
$\{i_1,i_2\}$ and $\{i_3,i_4\}$ agree.
For every $n>1$, the affine pair-deletion guarantee
\[
G=n\sum_i\ell_i-\binom n2\sum_i\ell_i^2
\]
satisfies
\[
G\le\frac{2}{3\sqrt3}\frac{n}{\sqrt n-1}
=\frac{2}{3\sqrt3}\sqrt n+O(1).
\]
Singer half-window templates asymptotically attain this coefficient.
\end{proposition}
\begin{proof}
First suppose $s\ge2$, and put $W=\sum_i\ell_i$.
The difference arcs $J_i-J_j$, indexed by ordered pairs $i\ne j$,
are pairwise disjoint. A common difference for $(i,j)$ and $(k,l)$
would give $x_i+x_l=x_k+x_j$; the label property and the two
off-diagonal conditions force $(i,j)=(k,l)$.
No difference arc has full measure, because it would intersect another
positive-length difference arc. Thus each has length $\ell_i+\ell_j$,
including the possibility of wraparound. Summing these lengths gives
\[
2(s-1)W\le1,\qquad s\le1+\frac1{2W}.
\]
By Cauchy--Schwarz,
\[
G\le nW-\frac{n(n-1)W^2}{2s}
\le nW-\frac{n(n-1)W^3}{1+2W}=:F(W).
\]
If $G\le0$, the assertion is immediate. Otherwise $F(W)>0$, so
$(n-1)W^2<1+2W$, or $W<1/(\sqrt n-1)$. Consequently
\[
\frac{n-1}{1+2W}\ge(\sqrt n-1)^2,
\qquad
G\le n\bigl(W-(\sqrt n-1)^2W^3\bigr).
\]
The maximum of $x-a^2x^3$ for $x\ge0$ is $2/(3\sqrt3 a)$;
apply this with $a=\sqrt n-1$. For $s=1$, directly maximizing
$n\ell-\binom n2\ell^2$ gives at most $n/(2(n-1))$, which is
also below the asserted bound.

For the final assertion, use arcs of length $1/(2m)$ based at the
$q+1$ elements of a Singer set divided by $m=q^2+q+1$.
Their label property follows from the same carry calculation as
Theorem~\ref{thm:compression}. At $q^2\sim3n/4$ they give the
coefficient of Theorem~\ref{thm:main}.
\end{proof}

The hypothesis allows different points carrying the same label in the
four-term equality. It is stronger than merely requiring that every
choice of one representative from each arc be Sidon; those statements
must not be interchanged. The proposition rules out improving the
coefficient solely by variable widths, non-grid centers, or mixtures
of these templates while retaining the same first--second moment
calculation. It leaves more accurate occupancy estimates open.

\begin{proposition}[Equal-length transversal templates]\label{prop:transversal-template-limit}
Let $J_1,\ldots,J_s\subset\T$ be disjoint arcs of a common positive
length $\ell$, with $s\ge5$. Suppose every choice of one point from
each arc is a Sidon set, with repeated summands allowed. For $n>1$,
the pair-deletion guarantee satisfies
\[
ns\ell-\binom n2s\ell^2
\le \frac{2}{3\sqrt3}\frac{n}{\sqrt n-1}
       +\frac{n}{n-1}
=\frac{2}{3\sqrt3}\sqrt n+O(1).
\]
Thus the coefficient in Proposition~\ref{prop:template-limit} remains
optimal for equal-length templates under this weaker hypothesis.
\end{proposition}
\begin{proof}
Discard arc endpoints, which changes neither lengths nor the displayed
guarantee. Write $c_i$ for the arc centers. Disjointness gives circular
separation at least $\ell$, and $s\ell\le1$, so $2\ell<1/2$.
There is at most one unordered pair of centers whose circular distance
is less than $2\ell$. Indeed, orient any such pair so that its positive
short difference $d$ lies in $[\ell,2\ell)$. Its oriented difference
arc is $(d-\ell,d+\ell)$ and contains $\ell$. Two disjoint close pairs
would therefore give equal differences from four distinct labels,
contradicting the transversal hypothesis. If two close pairs share a
center, lift that center to $0$ and their other centers to
$(-2\ell,2\ell)$. They cannot lie on the same side: their separation
would then be less than $\ell$. On opposite sides they are $-a,b$ with
$a,b\in[\ell,2\ell)$, and the point $(b-a)/2$ belongs to the central
arc. Together with the two outer centers it is a three-term arithmetic
progression, again a forbidden transversal. These real equalities also
hold on the circle, so wraparound causes no exception.

Delete the two arcs in the unique close pair, if one exists. All remaining
centers have circular separation at least $2\ell$. We claim that these
arcs satisfy the full label property of
Proposition~\ref{prop:template-limit}. First their oriented difference
arcs $J_i-J_j$, $i\ne j$, are pairwise disjoint. For pairs with the same
first index or the same second index, their centers are separated by at
least $2\ell$ and their open radii are $\ell$, proving disjointness.
For pairs with four distinct indices an intersection directly gives
a forbidden transversal collision. For pairs sharing an index in
opposite positions, an intersection gives, for example,
$x_i-x_j=x_k-x'_i$ modulo one. Average lifts of $x_i,x'_i$ in the same
open arc to obtain $z_i\in J_i$ and $2z_i=x_j+x_k$, a forbidden
three-term progression when $j\ne k$. For reversed pairs, average
the two representatives in each of the two arcs; this gives
$2z_i=2z_j$ with $z_i\ne z_j$, forbidden by the repeated-summand
Sidon condition. This exhausts all shared-index cases.

Every self-difference arc is $(-\ell,\ell)$, whereas center separation
at least $2\ell$ makes every off-diagonal difference arc disjoint from
it. Rearranging any four-term equality as equality of two differences
now proves the full label property: two self-differences give the same
unordered label multiset, a self-difference cannot equal an off-diagonal
one, and two off-diagonal differences must have the same ordered labels.
The use of lifts and open arcs also covers intervals crossing $0$.

Apply Proposition~\ref{prop:template-limit} to the remaining arcs.
Each of the at most two deleted arcs contributes at most
\[
\max_{x\ge0}\left(nx-\binom n2x^2\right)
=\frac{n}{2(n-1)}
\]
to the guarantee. Adding these contributions proves the bound.
Singer half-window templates supply asymptotic equality in the leading
coefficient, as in Proposition~\ref{prop:template-limit}.
\end{proof}

For fewer than five equal-length arcs the same guarantee is at most
$sn/(2(n-1))=O(1)$. This extension still bounds only the first--second
moment certificate, not actual occupancy or an adaptive choice of the
affine parameters. No assertion here extends the weaker transversal
hypothesis to arbitrary unequal lengths.

\subsection{Higher-point independence is false}

For the actual integer set $A=\{0,3,7\}$ and the arc $I=[0,1/6)$,
independent uniform $\theta,t$ satisfy
\begin{equation}\label{eq:negative-triple}
\Prob(t,t+3\theta,t+7\theta\in I)=\frac{13}{3024}
<\left(\frac16\right)^3=\frac{14}{3024}.
\end{equation}
Here and below circle membership is understood modulo one.
For a direct proof, put $l=1/6$, restrict $0\le t\le l$, and write
$u=t+3\theta-k$, $v=t+7\theta-j$ with $0\le u,v\le l$.
Then $7k-3j=4t-7u+3v\in[-7l,7l]$.
With $0\le\theta\le1$, the only possible integer pairs are
$(k,j)=(0,0),(3,7),(1,2),(2,5)$, apart from null boundaries.
The first two regions each have area $l^2/14=1/504$.
For $(1,2)$ the permitted $\theta$-interval has width $(8t-1)/42$
for $1/8\le t\le1/6$, giving area $1/6048$.
For $(2,5)$ its width is $1/126-4t/21$ on $0\le t\le1/24$,
giving the same area. The total is $13/3024$.

Inclusion--exclusion and two-point independence give union probability
$1273/3024$, below the independent-binomial value $1274/3024$.
This refutes exact independent-binomial domination. It does not refute
an asymptotic Poisson lower bound at the linear-modulus scale, and it
does not refute Erd\H{o}s's conjecture.

\subsection{Triple independence does not control heavy fibers}

The following example identifies a limitation of a proposed proof of
an improved occupancy bound from positive triple moments alone.

\begin{proposition}[Sharpness of pair deletion for general affine characters]
\label{prop:finite-field-occupancy}
Let $q$ be an odd prime, let $T\subseteq\F_q$ have size $1\le n\le q$,
and set $P_T=\{(x,x^2):x\in T\}\subseteq\F_q^2$. Choose
$u,v,t\in\F_q$ independently and uniformly, and let
\[
X=|\{x\in T:ux+vx^2+t=0\}|.
\]
Values of this affine map at any three distinct points are independent
and uniform. Nevertheless
\[
\Prob(X>0)=\frac n q-\frac{n(n-1)}{2q^2}
+\frac{(n-1)(n-2)}{2q^3}.
\]
In particular, if $n/q\to\lambda\in(0,1]$, this tends to
$\lambda-\lambda^2/2$, exactly the first--second-moment lower bound.
The set $P_T$ itself is Sidon.
\end{proposition}
\begin{proof}
For three distinct parameters $x_1,x_2,x_3$, the matrix with rows
$(1,x_i,x_i^2)$ has nonzero Vandermonde determinant. Hence the three
values are independent and uniform; the corresponding one- and
two-point assertions follow as well. Thus
$\E X=n/q$ and $\E\binom X2=\binom n2/q^2$.

Unless $u=v=t=0$, a nonzero polynomial of degree at most two has at
most two roots. The exceptional coefficient triple has probability
$q^{-3}$ and gives $X=n$. The identity
\[
\mathbf1_{r>0}=r-\binom r2+
\begin{cases}(r-1)(r-2)/2,&r\ge1,\\0,&r=0\end{cases}
\]
therefore gives the displayed exact formula. All the third factorial
moment, $\E\binom X3=\binom n3/q^3$, comes from this single rare
constant map, while its contribution to the occupancy surplus is
only $(n-1)(n-2)/(2q^3)=O(1/q)$.

Finally, if two pairs of parabola points have equal sums, their
parameters have the same sum $s$ and the same sum of squares $r$.
Since $q$ is odd, their product is $(s^2-r)/2$. Both unordered
parameter pairs are therefore the multiset of roots of the same
quadratic polynomial, including repeated roots. This proves the
Sidon assertion.
\end{proof}

At $\lambda=2/3$ this example gives limiting occupancy $4/9$ despite
exact three-point independence. It does not settle the corresponding
question for integer sets under continuous torus projection. It shows
that any argument using only character linearity and the first three
moments must either use the extra connected-group geometry or identify
a complementary class with a direct Sidon construction; this example
belongs to the latter class.

\subsection{An asymptotic Poisson obstruction at a fixed linear modulus}
\label{sec:grid-obstruction}

For a finite integer set $A$, write
\[
U_A(\epsilon)=\Prob_{\theta,t\in\T}
\bigl(\text{some }a\in A\text{ has }a\theta+t\in[0,\epsilon)\bigr),
\]
with independent uniform parameters. If $O_A(m)$ counts the occupied
half-cells $[j/m,(j+1/2)/m)$, translation invariance gives
$\E O_A(m)=mU_A(1/(2m))$.

\begin{theorem}[Failure of a universal Poisson occupancy estimate]
\label{thm:grid-poisson}
Put $d=32$, $\lambda=1/50$, and
\[
D=\frac1{12}\left(\frac23\right)^{32}\lambda^3>0.
\]
There are integer sets $A_L$ of size $n=L^{32}$ such that, with
$m=25n$,
\[
\limsup_{L\to\infty}U_{A_L}(1/(2m))
\le1-e^{-1/50}-\frac D2.
\]
The same conclusion holds along Singer moduli $m\sim25n$.
Thus the proposed uniform estimate
$U_A(1/(2m))\ge1-e^{-n/(2m)}-o(1)$ is false even when $m/n$
converges to a fixed positive constant. Here $D\approx1.54521\cdot10^{-12}$.
\end{theorem}

This concerns expected occupancy in the specified random projection.
It is not an upper bound on the largest Sidon subset, and says nothing
about failure of a Poisson estimate near its optimized parameter.

\begin{proof}
We give the counting details, since dependencies among five points matter.

\paragraph{Rational-plane counts.}
Let $2\le s<k\le5$, and let $W\subset\R^k$ be a rational $s$-plane
containing $\mathbf1=(1,\ldots,1)$. The lattice $K=W\cap\Z^k$ is
primitive. For an integral basis matrix $B$ of $K$, let $H(W)$ be the
maximum absolute value of an $s\times s$ minor. The gcd of these minors
is one: extend a basis of the primitive lattice to an integer basis of
$\Z^k$ and expand the resulting unimodular determinant. The signed
minors are the primitive Pl\"ucker coordinates of $W$; they determine
$W$ because $W=\{x:x\wedge p=0\}$ for its exterior vector $p$.
Consequently, with $J=\binom ks$, there are at most $(2H+1)^J$ planes
of height at most $H$, and hence also of height exactly $H$.

Projection onto coordinates attaining a largest minor maps $K$
injectively to a sublattice of $\Z^s$ of index $H$. A triangular
Hermite basis with positive diagonal entries $h_i$ has $\prod_i h_i=H$.
Counting coordinates successively, each in one congruence class, gives
\begin{equation}\label{eq:grid-box-count}
N_L(W):=|W\cap\{0,\ldots,L-1\}^k|
\le\prod_{i=1}^s(L/h_i+1)
\le L^s/H+(2^s-1)L^{s-1}.
\end{equation}

Consider ordered distinct $k$-tuples from $G_L=\{0,\ldots,L-1\}^d$
whose augmented vectors $(1,a_i)$ have rank $s$. The all-ones row and
the $d$ coordinate rows span a plane $W$ as above. It is not contained
in $x_i=x_j$, and a fixed plane permits at most $N_L(W)^d$ tuples.
Every realized plane has $H\le s!L^{s-1}$: wedge $\mathbf1$ with
$s-1$ independent coordinate rows. Its integer minors are bounded by
$s!L^{s-1}$ and form an integer multiple of the primitive Pl\"ucker
vector.

For any family of these planes with $H\ge h$, it follows that
\begin{equation}\label{eq:grid-height-sum}
\limsup_{L\to\infty}
\frac{\#\{\text{tuples in the family}\}}{L^{sd}}
\le\sum_W H(W)^{-d}
\le(h+1)(3h)^J h^{-d},
\end{equation}
provided $d\ge J+2$ and $d>(s-1)J$. These inequalities hold for $d=32$.
To justify the infinite sum, first restrict heights to $H\le R$ and
take $L\to\infty$ in \eqref{eq:grid-box-count}. For $H>R$, use
$(u+v)^d\le2^{d-1}(u^d+v^d)$. The normalized tail is at most
\[
2^{d-1}\sum_{H>R}O_k(H^J)H^{-d}
+O_{k,d}(L^{(s-1)J-d}).
\]
The second term vanishes as $L\to\infty$, and the first as $R\to\infty$.
Finally, $(2H+1)^J\le3^JH^J$ and the integral bound for a decreasing
power series give the final inequality in \eqref{eq:grid-height-sum}.
This order of limits keeps the factor $2^{d-1}$ out of the finite-height
main terms.

\paragraph{Height restrictions and the needed refinements.}
The following integral observation is useful. If $W$ has $k$ distinct
coordinate functions, then
\begin{equation}\label{eq:grid-packing-height}
k\le(H(W)+1)^{s-1}.
\end{equation}
Indeed, complete the primitive vector $\mathbf1$ to an integral basis
of $K$. The rows of its basis matrix are $(1,u_i)$, with distinct
$u_i\in\Z^{s-1}$. If two $u_i$ agree modulo $H+1$, their distinct
rows are linearly independent and extend with other rows to a nonzero
$s\times s$ minor divisible by $H+1$, a contradiction. Thus the $u_i$
have distinct residues modulo $H+1$.
In particular, rank-three five-tuples have height at least two.
For rank two, a basis $(\mathbf1,u)$ shows more explicitly that
$H=\max u_i-\min u_i\ge k-1$.

For collinear triples the primitive relation, up to sign, is a
permutation of $(r,H-r,-H)$ with $1\le r<H$ and $(r,H)=1$.
There are exactly $3\varphi(H)$ such planes. Their one-coordinate box
density is $1/H$: the middle coordinate is a convex combination of
the other two, with one divisibility condition of modulus $H$.
For fixed $H$ and $\epsilon<1/H$, the simultaneous three-hit probability
in $[0,\epsilon)$ is exactly $\epsilon^2/H$.
Indeed, write the three integer vectors as $a_0,a_0+u v,a_0+Hv$,
where $v\ne0$ is integral and $(u,H)=1$. Their phases are
$t',t'+u z,t'+Hz$, with $t'=t+a_0\cdot\theta$ and
$z=v\cdot\theta$ independent uniform circle points. Conditional on
the two endpoint phases, $z$ has $H$ equiprobable branches. Their
middle phases differ by the residues $u j/H$, which run through all
$H$ multiples of $1/H$. If the endpoints lie in $[0,\epsilon)$,
exactly one middle phase lies there, namely their real convex
combination. Fixed-height truncation and the tail argument above
justify summing this exact formula. Hence the collinear contribution
to the third factorial moment is bounded by $A_d\lambda^2$, where
\begin{equation}\label{eq:grid-Ad}
A_d=\frac14\,2^{-d}
+\frac12\left(1+\frac3{d-1}\right)3^{-d}.
\end{equation}
Here $(\lambda^2/2)\sum_{H\ge2}\varphi(H)H^{-d-1}$ is the preceding
bound; keep $H=2$ exactly and use $\varphi(H)\le H$ for the tail.

For rank-four five-tuples, a height-one relation has coefficients in
$\{-1,0,1\}$ summing to zero. Support two would equate two points.
The only remaining case has two coefficients $+1$, two $-1$, and one
zero. There are 15 such planes. Their one-coordinate count is
$L(2L^3+L)/3$, so each has density $2/3$. The other rank-four planes,
of height at least two, have normalized count at most
$23328\,2^{-d}$ by \eqref{eq:grid-height-sum}.

There are at most 2100 height-two rank-three planes with five distinct
coordinate functions. Their primitive maximal minors have gcd one
and absolute value at most two, so some minor is $\pm1$. Normalize
the corresponding three indexed affine points to $(0,0),(1,0),(0,1)$.
Each remaining integer point $(x,y)$ satisfies
$|x|,|y|,|1-x-y|\le2$. Excluding the three frame points leaves 15
candidates: there are respectively $2,3,2,4,4$ at
$x=-2,-1,0,1,2$. There are at most $\binom53\cdot15\cdot14=2100$
choices; ignoring other triangle minors only enlarges this count.
For heights at least three, summing $(2H+1)^{10}H^{-d}$ gives the bound
$3^{20}(1+3/(d-11))3^{-d}$. Therefore the normalized rank-three
five-tuple count is at most
\begin{equation}\label{eq:grid-Bd}
B_d=2100\,2^{-d}
+3^{20}\left(1+\frac3{d-11}\right)3^{-d}.
\end{equation}

For rank-two five-tuples, normalize the primitive integer positions
to minimum zero and maximum $H\ge4$. After identifying reflection,
there are at most ten choices of the positions of the extrema and
at most $H^3$ ordered choices for the other positions. Thus at most
$10H^3$ planes have height $H$, giving normalized count at most
\begin{equation}\label{eq:grid-Cd}
C_d=640\left(1+\frac4{d-4}\right)4^{-d}.
\end{equation}
All these refinements use the same tail justification as
\eqref{eq:grid-height-sum}.

\paragraph{The moment comparison.}
Fix $d=32$, take $n=L^d$ and $\epsilon=\lambda/n$, and let
\[
X_L=\sum_{a\in G_L}\mathbf1_{[0,\epsilon)}(t+a\cdot\theta),
\]
where $t,\theta_1,\ldots,\theta_d$ are independent uniform circle
points. A tuple of augmented rank $s$ has $s$ independent uniform
phase coordinates: its full-row-rank integral matrix defines a
surjective torus homomorphism. Consequently its simultaneous hit
probability is at most $\epsilon^s$, and is exactly $\epsilon^k$
for a full-rank $k$-tuple. With $M_j=\E\binom{X_L}{j}$, we obtain
\[
M_1=\lambda,\qquad M_2\longrightarrow\lambda^2/2,\qquad
\limsup M_3\le\lambda^3/6+A_d\lambda^2.
\]

Almost all ordered quadruples have full augmented rank four: the
lower-rank counts are $O_d(n^3)$. Their fourth-moment contribution
tends to $\lambda^4/24$. The grid's ordered additive energy is
$[(2L^3+L)/3]^d$. Repeated-index solutions contribute $O(n^2)$,
and collinear quadruples contribute $O_d(n^2)$ by the rank-two count.
Thus the number of unordered noncollinear parallelograms is
$(1/8)(2/3)^d n^3+o(n^3)$. Each has simultaneous hit probability
$2\epsilon^3/3$: three independent phases determine the fourth, and
the integral is the square-integral of the triangular convolution of
$\mathbf1_{[0,\epsilon)}$ with itself. It follows that
\[
\liminf M_4\ge\lambda^4/24+
\frac1{12}(2/3)^d\lambda^3.
\]
Separating the fifth moment by augmented ranks five, four, three,
and two gives
\[
\limsup M_5\le\frac{\lambda^5}{120}
+\frac18(2/3)^d\lambda^4
+\frac{23328}{120}\,2^{-d}\lambda^4
+\frac{B_d\lambda^3+C_d\lambda^2}{120}.
\]
The divisors $6,24,120$ throughout convert ordered distinct tuples
to subsets.

The fifth Bonferroni inequality gives
$\Prob(X_L>0)\le M_1-M_2+M_3-M_4+M_5$.
For $0<\lambda\le1$, the degree-five alternating exponential
polynomial exceeds $1-e^{-\lambda}$ by at most $\lambda^6/720$.
Consequently
\[
\limsup\Prob(X_L>0)\le1-e^{-\lambda}-D_d+E_d,
\]
where
\[
D_d=\frac1{12}(2/3)^d\lambda^3
\]
and
\[
E_d=\frac{\lambda^6}{720}+A_d\lambda^2
+\frac18(2/3)^d\lambda^4
+\frac{23328}{120}\,2^{-d}\lambda^4
+\frac{B_d\lambda^3+C_d\lambda^2}{120}.
\]
At $d=32$, $\lambda=1/50$, substitution into these explicit fractions gives
\[
\frac{E_d}{D_d}
=\frac{290086345369360189976620730533563}
{1311684514281872654556200960000000}
<\frac12.
\]
The script \texttt{research/sidon/scripts/poisson\_grid\_constants.py}
reproduces the rational inequality without floating-point comparisons.
This proves the claimed deficit in the grid model.

\paragraph{Integer realization and Singer moduli.}
For each $L$, choose a sufficiently large integer base $M>2(L-1)$
and form
\[
A_{L,M}=\left\{\sum_{j=0}^{d-1}a_jM^j:0\le a_j<L\right\}.
\]
It has $L^d$ distinct integer elements. For fixed $L,d$, the
distribution of $(t,\theta,M\theta,\ldots,M^{d-1}\theta)$ converges
to Haar measure on $\T^{d+1}$ as $M\to\infty$: a fixed nontrivial
Fourier character has zero expectation for all sufficiently large
$M$, since a nonzero integer polynomial has finitely many roots.
Approximation by trigonometric polynomials gives weak convergence.
The boundary of the union-hit event is contained in finitely many
affine subtori and has measure zero. Choose $M=M(L)$ so the resulting
probability differs from that of the grid by at most $1/L$.
This proves the integer claim at $m=25n$, for which
$1/(2m)=\lambda/n$ exactly.

Alternatively choose primes $q\sim5\sqrt n$ and
$m=q^2+q+1\sim25n$. Then $n/(2m)\to\lambda$. Changing interval
length from $\lambda/n$ to $1/(2m)$ changes any union probability
by at most $n|1/(2m)-\lambda/n|=o(1)$, by the union bound.
The same positive deficit survives, including the integer realization.
\end{proof}

These integer digit grids also satisfy Corollary~\ref{cor:box-model},
since $K=L-1=o(\sqrt{L^{32}})$. They therefore have the stronger
$0.544331\ldots$ Sidon lower bound proved for that class. This makes
the distinction between an occupancy obstruction and a Sidon
obstruction particularly explicit.

\subsection{A small quotient can conceal an easy deletion}

Let $A_k=\{0,1,2,4,\ldots,2^k\}$, so $n=k+2$. Its only nontrivial
unordered collision equalities are
\[
0+2^{j+1}=2^j+2^j\qquad(0\le j<k).
\]
This follows from binary expansion: two distinct positive powers give
two nonzero digits, while a doubled power and a sum with zero give
one. Each displayed relation eliminates the next power with a unit
coefficient, leaving the universal Freiman quotient free on the
generators at $0$ and $1$. It therefore has rank two. Nonetheless,
removing $0$ leaves a Sidon set of size $n-1$.
Its ordered additive energy is only
\[
E(A_k)=2n^2-n+4(n-2)=2n^2+3n-8:
\]
each displayed collision merges ordered multiplicities $1$ and $2$,
increasing energy by $3^2-1^2-2^2=4$.
Thus neither low collision energy nor a small universal quotient
automatically supplies higher character independence or identifies
the difficult branch of Sidon extraction.

\subsection{A precise high-energy reduction}
\label{sec:energy-reduction}

The following is a standard probabilistic alteration argument, included
to isolate a concrete remaining branch. It is not claimed as a novel
universal improvement.

Let $r_A(s)=|\{(a,b)\in A^2:a+b=s\}|$ and let
$E(A)=\sum_s r_A(s)^2$. Let $T_3(A)$ count unordered three-element
arithmetic progressions in $A$, and let $Q_4(A)$ count unordered pairs
of distinct off-diagonal two-element subsets of $A$ having the same
sum. Such a pair necessarily uses four distinct elements.

\begin{proposition}[Energy-dependent alteration]
\label{prop:energy-alteration}
For every finite integer set $A$ of size $n$,
\[
s(A)\ge\ceil{\max_{0\le p\le1}
\bigl(np-T_3(A)p^3-Q_4(A)p^4\bigr)}.
\]
Moreover,
\[
E(A)=2n^2-n+4T_3(A)+8Q_4(A),\qquad
T_3(A)\le\frac{\sqrt{nE(A)}-n}{2}.
\]
For every fixed $\gamma>0$, uniformly over sets satisfying
$E(A)\le\gamma n^{5/2}$,
\[
s(A)\ge
\frac34\left(\frac2\gamma\right)^{1/3}\sqrt n
-\gamma^{-1/2}n^{1/4}
\]
for all sufficiently large $n$.
\end{proposition}
\begin{proof}
Select every element independently with probability $p$. Every
nontrivial Sidon relation uses either three distinct elements, forming
an arithmetic progression, or four distinct elements counted by
$Q_4$. Delete one point from each surviving bad relation, in any
order. Deletion creates no relation. The remaining set is Sidon and
has size at least the selected cardinality minus the original number
of surviving bad relations. Taking expectations gives the finite
bound, including its ceiling and optimization over $p$.

For the energy identity, a sum with $k$ unordered off-diagonal
representations and $e\in\{0,1\}$ diagonal representations has
ordered representation count $2k+e$. Its contribution beyond the
trivial ordered equalities is
\[
(2k+e)^2-(4k+e)=8\binom k2+4ke.
\]
The first term counts proper four-point collisions and the second
counts arithmetic progressions. Summing proves the identity.
Also
\[
T_3(A)=\frac12\left(\sum_{b\in A}r_A(2b)-n\right)
\le\frac{\sqrt{nE(A)}-n}{2},
\]
by Cauchy--Schwarz. In particular, diagonal relations cannot simply be
discarded using only the weaker bound $T_3\le n^2$.

Suppose now that $E(A)\le\gamma n^{5/2}$. Put
$c=(2/\gamma)^{1/3}$ and $p=c/\sqrt n\le1$. The fourth-order loss is
at most $(\gamma c^4/8)\sqrt n$, since $Q_4\le E(A)/8$.
The third-order loss is at most
\[
\frac{c^3}{n^{3/2}}\frac{\sqrt\gamma}{2}n^{7/4}
=\gamma^{-1/2}n^{1/4}.
\]
Finally $c-\gamma c^4/8=3c/4$, proving the stated bound.
\end{proof}

Consequently the constant $0.4$ is already supplied by this argument
on the branch
\[
E(A)\le\frac{3375}{256}\,n^{5/2},
\]
up to its vanishing asymptotic loss. The constant $0.5$ is supplied on
the branch $E(A)\le(27/4)n^{5/2}$. A universal improvement to either
constant would still require a proof for the complementary high-energy
branch. No such complementary proof is established here.

\section{Finite verification and reproducibility}

The independent Python script
\texttt{research/sidon/computation/verify\_affine\_sidon.py}
uses exact integer and rational arithmetic and no third-party packages.
It constructs Singer sets for primes $q=2,3,5,7,11,13,17,19$ by explicit
cubic finite fields, verifies every nonzero cyclic difference occurs
once, and checks all unordered sums including repeated summands.
It exhausts a prime-grid version of the two-point calculation, and
tests exact rational affine maps on intervals, negative intervals,
separated clusters, geometric progressions, Boolean cubes, and dense
and sparse random sets. Every original two-sum equality in every
retained test domain is checked before every Singer pullback is checked.

The recorded run uses seed $20261003$ and $2000$ affine trials for each
of seven ambient examples. Its output is
\texttt{research/sidon/computation/verification.json}, including source
sets, explicit Sidon subsets, affine parameters, target translates,
and a SHA-256 hash of the script. Reproduce it from the project root by
\begin{verbatim}
python3 research/sidon/computation/verify_affine_sidon.py --trials 2000
\end{verbatim}
The widened-window counterexample is verified exactly with
$m=5$, $\theta=13/40$, $t=3/40$, $\delta=3/4$.
The small negative-correlation example also has a separate rational
certificate, but its area proof above is independent of computation.

Further exact certificates are provided by the standard-library scripts
\path{research/sidon/scripts/ap_surplus_certificate.py},
\path{research/sidon/scripts/poisson_grid_constants.py},
\path{research/sidon/scripts/rigidifier_certificate.py}, and
\path{research/sidon/scripts/short_window_certificate.py}.
They check, respectively, the progression surplus and representative
probabilities; the printed rational grid deficit; and the integral
relations and collapse properties of the two rigidifier constructions.
Their saved JSON files are in \path{research/sidon/computation/}
and \path{research/sidon/results/}.
The progression certificate exhausts all $512$ subsets of a nine-point
interval and $15\,870$ exact finite sampling parameters.
These certificates support fully written proofs rather than replace them.

An optional Z3 search, \path{research/sidon/scripts/cluster_sidon_search.py},
records explicit finite Sidon witnesses in
\path{research/sidon/results/cluster_sidon_search.json}. Each witness is
independently checked by enumerating all unordered two-sums. Solver
UNSAT results are diagnostic unless accompanied by a separately checked
proof certificate; no asymptotic conclusion depends on them.

Exploratory searches for substantially stronger cyclic models did not
yield a universal theorem or a general counterexample. In particular, solver timeouts were
\emph{UNKNOWN}, not UNSAT, and random-subset search failures were not
interpreted as negative mathematical evidence.

\section{Remaining claims and session record}

The conjectural constant $1$ remains open in this report. The next
mathematical steps are limited to the following three explicit claims;
none is proved here.
\begin{enumerate}
\item At $m\sim3n/4$, prove a universal occupied-arc surplus
\[
U_A(1/(2m))\ge\frac49+\eta-o(1)
\]
for some fixed $\eta>0$. Since expected occupied cells equal $mU_A$,
this would increase the Sidon coefficient by $(\sqrt3/2)\eta$.
The Poisson obstruction at $m\sim25n$ does not disprove this weaker
claim at the indicated parameter.
\item Prove a structural dichotomy that either supplies a subset of
size $(\alpha+o(1))n$, for some fixed $\alpha>1/2$, with distinct
classes in its rebuilt modular Freiman quotient at primes
$p\sim\sqrt{3\alpha n/2}$, or gives a direct Sidon bound exceeding
$2/(3\sqrt3)$ in the complementary case. The modular branch would
supply coefficient $(2\sqrt2/(3\sqrt3))\sqrt\alpha$.
The rigidifier and robust-product theorems explain why near-lossless
finite modeling cannot simply be assumed.
\item Prove $s(A)\ge(0.4-o(1))\sqrt n$ in the remaining regime
\[
E(A)>\frac{3375}{256}n^{5/2},\qquad
T(A)<\frac{2}{27}n^2.
\]
The energy alteration and progression-sensitive compression already
supply the same asymptotic coefficient on the other two branches.
A proof of this precise residual case would therefore yield a
universal $0.4$ bound. No conclusion for this case follows from global
energy or a small quotient alone.
\end{enumerate}

The research session ran from 22:30:34 UTC on 3 October 2026 to
00:30:51 UTC on 4 October 2026: 2 hours, 0 minutes, 17 seconds.
Independent internal adversarial audits found no gap in the final
claims; these are not external peer review or publication-priority
certification. The source, exact certificates, and detailed audit
records accompany this report.

\clearpage
\begin{thebibliography}{9}
\bibitem{BR}
A. Bailleul and R. Riblet,
\emph{On the largest Sidon subset in a finite subset of $\R^N$},
arXiv:2605.03181v1 (4 May 2026), especially Lemmas 2.3--2.5 and
Theorem 2.2. \url{https://arxiv.org/html/2605.03181v1}.

\bibitem{RuzsaII}
I. Z. Ruzsa, \emph{Solving a linear equation in a set of integers II},
Acta Arith. \textbf{72} (1995), 385--397, Theorem 4.4.
\url{https://matwbn.icm.edu.pl/ksiazki/aa/aa72/aa7248.pdf}.

\bibitem{Ruzsa2006}
I. Z. Ruzsa, \emph{Additive and multiplicative Sidon sets},
Acta Math. Hungar. \textbf{112} (2006), 345--354.
\url{https://doi.org/10.1007/s10474-006-0102-0}.
Full text not obtained; no theorem in this report relies on it.

\bibitem{Singer}
J. Singer, \emph{A theorem in finite projective geometry and some
applications to number theory}, Trans. Amer. Math. Soc.
\textbf{43} (1938), 377--385.
\url{https://doi.org/10.1090/S0002-9947-1938-1501951-4}.

\bibitem{PZ}
J. Pach and D. Zakharov, \emph{Ruzsa's problem on bi-Sidon sets},
Combinatorica \textbf{45} (2025), article 26, Propositions 1--2.
\url{https://arxiv.org/html/2409.03128v2}.

\bibitem{Huber}
F. Huber, \emph{Large Sidon subsets of distinct multinomial
coefficients}, arXiv:2608.02667v1 (2 August 2026).
\url{https://arxiv.org/html/2608.02667v1}.

\bibitem{RibletIntervals}
R. Riblet, \emph{Sidon sets in a union of intervals},
arXiv:2202.01296, Theorem 3.1(ii).
\url{https://arxiv.org/pdf/2202.01296}.

\bibitem{Green}
B. Green, \emph{Additive combinatorics, exercises 2}, Question 7.
\url{https://people.maths.ox.ac.uk/greenbj/papers/addcomb-sheet2.pdf}.
The relevant assertion is recorded as unverified in this report.
\end{thebibliography}
\end{document}
